% Zumdahl Chapter 5 comprehensive test -- 105 free-response items, built to a
% teacher request covering all 20 Interactive Examples (5.1-5.20), 72 named
% end-of-chapter exercises, and all 13 Challenge Problems (Q157-Q169). Six
% sessions: Part A-B (examples), Part C (gas-law exercises), Part D
% (density/stoichiometry exercises), Part E (partial-pressure/atmosphere
% exercises), Part F (KMT/real-gas exercises), Part G (challenge problems).
% Every numeric answer independently recomputed in Python before being typed
% here, at full precision, rounding only at the end.
%
% Items adapted from Zumdahl Ch5 examples/exercises are reworded as original
% items with the source numbers preserved and tagged "after Zumdahl Ex./Qn"
% -- the public repo carries no Cengage text (convention 6).
\documentclass{apchem-exam}

% Same fix as Units 4/5/6: siunitx groups any run of 5+ digits, and this
% paper carries 5+ digit Avogadro-scale counts and molecule concentrations.
\sisetup{group-minimum-digits = 7}

\apcunitword{Chapter}
\apcunit{5}
\apcunittitle{Gases}
\apcdoctitle{Comprehensive Test}
\apcsubtitle{Zumdahl \S5.1--5.9 \textbullet\ 105 questions, 376 points,
  6 sessions \textbullet\ S1: Parts A--B (Examples), 20Q/72pts/90min
  \textbullet\ S2: Part C (Gas Laws), 24Q/73pts/90min \textbullet\ S3:
  Part D (Density \& Stoichiometry), 16Q/54pts/70min \textbullet\ S4:
  Part E (Partial Pressure \& Atmosphere), 19Q/63pts/80min \textbullet\
  S5: Part F (KMT \& Real Gases), 13Q/39pts/50min \textbullet\ S6: Part G
  (Challenge), 13Q/75pts/100min}

\begin{document}
\apctitleblock

\begin{apcnote}[How this paper runs]
\textbf{Six sessions, not one.} Session 1 (Parts A--B) works through all
twenty of the textbook's own worked examples, rebuilt as original problems
using the same numbers and setups, so you can check your method directly
against the book. Sessions 2--5 (Parts C--F) are the full end-of-chapter
exercise set, grouped by topic. Session 6 (Part G) is the Challenge Problem
set -- the hardest material in the chapter. Do not attempt more than one
session in a single sitting.

Calculator, periodic table, and equations sheet are permitted throughout.
\textbf{Show every step.} A bare final answer earns the answer mark only;
the setup marks are for work you can point to. Carry units through each
line and round only at the end. Use $R = \SI{0.08206}{\liter\atm\per
\mole\per\kelvin}$ unless a part specifically calls for SI units.
\end{apcnote}

\begin{apcnote}[Scope --- read before Session 5]
Three tools this chapter uses are \textbf{not} on the AP equation sheet and
not part of CED Topics 3.4--3.6: the \textbf{root-mean-square speed}
formula ($u_{\text{rms}} = \sqrt{3RT/M}$), \textbf{Graham's law of
effusion} ($\text{rate}_1/\text{rate}_2 = \sqrt{M_2/M_1}$), and the
\textbf{van der Waals equation} itself as a numerical tool. The CED gives
you $KE = \tfrac12 mv^2$ and the qualitative idea that Kelvin temperature
tracks average kinetic energy (3.5.A.2--3.5.A.3), and it gives you the
qualitative reasoning that interparticle attraction and particle volume
explain non-ideal behavior (3.6.A.1) -- but neither an rms-speed equation,
an effusion-rate equation, nor van der Waals constants are anywhere in the
official equation sheet. Every item using one of these three tools is
flagged \textbf{[beyond CED]}. It is here because this is standard
introductory chemistry and because reading a molecular-speed or
non-ideal-gas argument qualitatively is core CED content even when solving
for a number is not.
\end{apcnote}

% ============================================================================
% SESSION 1
% ============================================================================
\examsection{A}{Worked-Example Mastery I --- Pressure and the Empirical Gas
  Laws}{10 questions \textbullet\ 37 points}{\calculatorok}

% ---------------------------------------------------------------------------
\frq{short}{3}{Zumdahl \S5.1 \textbullet\ pressure unit conversion
  \emph{(after Zumdahl Example 5.1)}}

A student's gas sample is measured at \SI{49}{\torr}.

Convert this pressure to (a) atmospheres and (b) pascals.
\workspace[3cm]
\keyonly{\begin{apcanswer}\textbf{(a)} $49 \div 760 = \mathbf{0.064}$~atm

\textbf{(b)} $0.064 \times \SI{101325}{\pascal\per\atm}
= \mathbf{6.5\times10^3}$~Pa\end{apcanswer}}

\begin{rubric}
  \rpoint{1}{divides by 760 to reach \SI{0.064}{\atm}}
  \rpoint{1}{converts atm to Pa using \SI{101325}{\pascal\per\atm}}
  \rpoint{1}{correct sig figs (2) and units on both answers}
\end{rubric}

% ---------------------------------------------------------------------------
\frq{short}{3}{Zumdahl \S5.2 \textbullet\ Boyle's law
  \emph{(after Zumdahl Example 5.2)}}

A \SI{1.53}{\liter} sample of gaseous \ce{SO2} at \SI{5.6e3}{\pascal} has
its pressure changed to \SI{1.5e4}{\pascal} at constant temperature.

Calculate the new volume.
\workspace[3.5cm]
\keyonly{\begin{apcanswer}\textbf{Boyle's law:} $P_1V_1 = P_2V_2$

\[V_2 = \frac{P_1V_1}{P_2} = \frac{(5.6\times10^3)(1.53)}{1.5\times10^4}
= \mathbf{0.57}\ \text{L}\]\end{apcanswer}}

\begin{rubric}
  \rpoint{1}{sets up $P_1V_1=P_2V_2$}
  \rpoint{1}{correct algebraic rearrangement for $V_2$}
  \rpoint{1}{\SI{0.57}{\liter}, correct sig figs}
\end{rubric}

% ---------------------------------------------------------------------------
\frq{long}{7}{Zumdahl \S5.2 \textbullet\ Boyle's law and the approach to
  ideal behavior \emph{(after Zumdahl Example 5.3)}}

A sample of \SI{1.0}{\mole} \ce{NH3} gas at \SI{0}{\celsius} was studied at
six pressures, and its volume was measured at each:

\begin{center}
\begin{tabular}{cc}
\toprule
$P$ (atm) & $V$ (L) \\
\midrule
0.1300 & 172.1 \\
0.2500 & 89.28 \\
0.3000 & 74.35 \\
0.5000 & 44.49 \\
0.7500 & 29.55 \\
1.000 & 22.08 \\
\bottomrule
\end{tabular}
\end{center}

Calculate $PV$ at each pressure, then state what value $PV$ appears to
approach as $P \to 0$, and explain what that value represents physically.
\workspace[7cm]
\keyonly{\begin{apcanswer}
\[
\begin{array}{cc}
P\ (\text{atm}) & PV\ (\text{L}\cdot\text{atm})\\
0.1300 & 22.373\\
0.2500 & 22.320\\
0.3000 & 22.305\\
0.5000 & 22.245\\
0.7500 & 22.163\\
1.000 & 22.080
\end{array}
\]

As pressure drops toward zero, $PV$ climbs toward roughly
\textbf{22.4~L$\cdot$atm}, essentially $nRT=(1.0)(0.08206)(273)
=22.4$~L$\cdot$atm. This is the \textbf{ideal-gas limit}: every real gas
behaves more like an ideal gas at low pressure, because the molecules
spend more time far apart, where intermolecular attractions and their
own volume matter least.
\end{apcanswer}}

\begin{rubric}
  \rpoint{1}{$PV=22.373$ at $P=0.1300$}
  \rpoint{1}{$PV=22.320$ at $P=0.2500$}
  \rpoint{1}{$PV=22.305$ at $P=0.3000$}
  \rpoint{1}{$PV=22.245$ at $P=0.5000$}
  \rpoint{1}{$PV=22.163$ at $P=0.7500$}
  \rpoint{1}{$PV=22.080$ at $P=1.000$}
  \rpoint{1}{identifies the low-pressure limit as $nRT\approx\SI{22.4}
    {\liter\atm}$ and ties it to ideal behavior}
\end{rubric}

% ---------------------------------------------------------------------------
\frq{short}{3}{Zumdahl \S5.3 \textbullet\ Charles's law
  \emph{(after Zumdahl Example 5.4)}}

A gas sample has a volume of \SI{2.58}{\liter} at \SI{15}{\celsius} and
some fixed pressure.

Find its volume at \SI{38}{\celsius}, same pressure.
\workspace[3.5cm]
\keyonly{\begin{apcanswer}$T_1 = 15+273 = \SI{288}{\kelvin}$,
$T_2 = 38+273 = \SI{311}{\kelvin}$

\[V_2 = V_1 \times \frac{T_2}{T_1} = 2.58 \times \frac{311}{288}
= \mathbf{2.79}\ \text{L}\]\end{apcanswer}}

\begin{rubric}
  \rpoint{1}{converts both temperatures to kelvin}
  \rpoint{1}{sets up Charles's law $V_2=V_1(T_2/T_1)$}
  \rpoint{1}{\SI{2.79}{\liter}}
\end{rubric}

% ---------------------------------------------------------------------------
\frq{short}{3}{Zumdahl \S5.3 \textbullet\ Avogadro's law
  \emph{(after Zumdahl Example 5.5)}}

A \SI{12.2}{\liter} sample at \SI{1}{\atm}, \SI{25}{\celsius} contains
\SI{0.50}{\mole} \ce{O2}. Every \ce{O2} molecule present is converted to
ozone via \ce{3O2(g) -> 2O3(g)}, at the same temperature and pressure.

Find the volume of ozone gas produced.
\workspace[3.5cm]
\keyonly{\begin{apcanswer}\textbf{Step 1 --- moles of \ce{O3}} from the 3:2
ratio:

\[n(\ce{O3}) = 0.50 \times \frac{2}{3} = 0.333\ \text{mol}\]

\textbf{Step 2 --- Avogadro's law} ($V\propto n$ at fixed $T,P$):

\[V_2 = V_1 \times \frac{n_2}{n_1} = 12.2 \times \frac{0.333}{0.50}
= \mathbf{8.1}\ \text{L}\]\end{apcanswer}}

\begin{rubric}
  \rpoint{1}{finds \SI{0.33}{\mole} \ce{O3} from the 3:2 stoichiometric
    ratio}
  \rpoint{1}{applies Avogadro's law, $V\propto n$}
  \rpoint{1}{\SI{8.1}{\liter}}
\end{rubric}

% ---------------------------------------------------------------------------
\frq{short}{3}{Zumdahl \S5.4 \textbullet\ ideal gas law, solve for moles
  \emph{(after Zumdahl Example 5.6)}}

A sample of \ce{H2} gas occupies \SI{8.56}{\liter} at \SI{0}{\celsius} and
\SI{1.5}{\atm}.

Find the moles of \ce{H2} present.
\workspace[3.5cm]
\keyonly{\begin{apcanswer}\[n = \frac{PV}{RT} = \frac{(1.5)(8.56)}
{(0.08206)(273)} = \mathbf{0.57}\ \text{mol}\]\end{apcanswer}}

\begin{rubric}
  \rpoint{1}{converts to \SI{273}{\kelvin}}
  \rpoint{1}{rearranges $PV=nRT$ correctly for $n$}
  \rpoint{1}{\SI{0.57}{\mole}}
\end{rubric}

% ---------------------------------------------------------------------------
\frq{short}{3}{Zumdahl \S5.4 \textbullet\ ideal gas law, solve for pressure
  \emph{(after Zumdahl Example 5.7)}}

Ammonia gas held at \SI{1.68}{\atm} in a \SI{7.0}{\milli\liter} container
is compressed to \SI{2.7}{\milli\liter} at constant temperature.

Find the final pressure.
\workspace[3.5cm]
\keyonly{\begin{apcanswer}Constant $n,T$ means Boyle's law applies:

\[P_2 = P_1 \times \frac{V_1}{V_2} = 1.68 \times \frac{7.0}{2.7}
= \mathbf{4.4}\ \text{atm}\]\end{apcanswer}}

\begin{rubric}
  \rpoint{1}{recognizes constant $n,T$ allows Boyle's law}
  \rpoint{1}{correct algebra}
  \rpoint{1}{\SI{4.4}{\atm}}
\end{rubric}

% ---------------------------------------------------------------------------
\frq{short}{3}{Zumdahl \S5.4 \textbullet\ ideal gas law, solve for volume
  \emph{(after Zumdahl Example 5.8)}}

Methane gas occupies \SI{3.8}{\liter} at \SI{5}{\celsius}. It is heated at
constant pressure to \SI{86}{\celsius}.

Find the new volume.
\workspace[3.5cm]
\keyonly{\begin{apcanswer}$T_1=\SI{278}{\kelvin}$, $T_2=\SI{359}{\kelvin}$

\[V_2 = V_1\times\frac{T_2}{T_1} = 3.8\times\frac{359}{278}
= \mathbf{4.9}\ \text{L}\]\end{apcanswer}}

\begin{rubric}
  \rpoint{1}{converts both temperatures to kelvin}
  \rpoint{1}{applies Charles's law correctly}
  \rpoint{1}{\SI{4.9}{\liter}}
\end{rubric}

% ---------------------------------------------------------------------------
\frq{short}{4}{Zumdahl \S5.4 \textbullet\ combined gas law
  \emph{(after Zumdahl Example 5.9)}}

A sample of diborane, \ce{B2H6(g)}, is at \SI{345}{\torr},
\SI{-15}{\celsius}, \SI{3.48}{\liter}. Conditions change to
\SI{36}{\celsius} and \SI{468}{\torr}.

Find the new volume. (Pressure may stay in torr throughout --- the unit
cancels in the ratio.)
\workspace[4cm]
\keyonly{\begin{apcanswer}$T_1 = \SI{258}{\kelvin}$, $T_2=\SI{309}{\kelvin}$

\[V_2 = V_1\times\frac{P_1}{P_2}\times\frac{T_2}{T_1}
= 3.48\times\frac{345}{468}\times\frac{309}{258}
= \mathbf{3.07}\ \text{L}\]\end{apcanswer}}

\begin{rubric}
  \rpoint{1}{converts both temperatures to kelvin}
  \rpoint{1}{sets up the combined gas law correctly}
  \rpoint{1}{recognizes torr does not need conversion here}
  \rpoint{1}{\SI{3.07}{\liter}}
\end{rubric}

% ---------------------------------------------------------------------------
\frq{short}{5}{Zumdahl \S5.4 \textbullet\ combined gas law with unit
  conversion \emph{(after Zumdahl Example 5.10)}}

\SI{0.35}{\mole} argon gas is held at \SI{13}{\celsius} and \SI{568}{\torr}.
Conditions change to \SI{56}{\celsius} and \SI{897}{\torr}.

Find the change in volume, $V_2-V_1$.
\workspace[6cm]
\keyonly{\begin{apcanswer}\textbf{Convert everything first.}
$P_1 = 568/760=\SI{0.7474}{\atm}$, $P_2=897/760=\SI{1.1803}{\atm}$,
$T_1=\SI{286}{\kelvin}$, $T_2=\SI{329}{\kelvin}$.

\textbf{Solve each state separately} (moles are fixed, but pressure and
temperature both change, so there is no single-ratio shortcut for
$\Delta V$):

\[V_1 = \frac{nRT_1}{P_1} = \frac{(0.35)(0.08206)(286)}{0.7474}
= 11.0\ \text{L}\]

\[V_2 = \frac{nRT_2}{P_2} = \frac{(0.35)(0.08206)(329)}{1.1803}
= 8.01\ \text{L}\]

\[\Delta V = V_2-V_1 = 8.01-11.0 = \mathbf{-3.0}\ \text{L}\]

The volume \textbf{decreases} by about \SI{3.0}{\liter} --- the pressure
rose by a larger factor than the temperature did, so compression wins out
over thermal expansion.\end{apcanswer}}

\begin{rubric}
  \rpoint{1}{converts both pressures to atm}
  \rpoint{1}{converts both temperatures to kelvin}
  \rpoint{1}{finds $V_1\approx\SI{11.0}{\liter}$}
  \rpoint{1}{finds $V_2\approx\SI{8.01}{\liter}$}
  \rpoint{1}{$\Delta V\approx\SI{-3.0}{\liter}$, with the sign/direction
    interpreted}
\end{rubric}

% ============================================================================
\clearpage
\examsection{B}{Worked-Example Mastery II --- Stoichiometry, Density,
  Dalton's Law, KMT}{10 questions \textbullet\ 35 points}{\calculatorok}

% ---------------------------------------------------------------------------
\frq{short}{3}{Zumdahl \S5.4 \textbullet\ gas stoichiometry, STP molar
  volume \emph{(after Zumdahl Example 5.11)}}

A sample of \ce{N2} gas occupies \SI{1.75}{\liter} at STP.

Find the moles of \ce{N2} present, using the molar volume of a gas at STP.
\workspace[3cm]
\keyonly{\begin{apcanswer}\[n = \frac{1.75\ \text{L}}{22.42\ \text{L/mol}}
= \mathbf{0.0781}\ \text{mol}\]\end{apcanswer}}

\begin{rubric}
  \rpoint{1}{identifies that the STP molar-volume shortcut applies}
  \rpoint{1}{correct division}
  \rpoint{1}{\SI{0.0781}{\mole}, correct sig figs}
\end{rubric}

% ---------------------------------------------------------------------------
\frq{short}{4}{Zumdahl \S5.4 \textbullet\ gas stoichiometry, decomposition
  \emph{(after Zumdahl Example 5.12)}}

\ce{CaCO3(s) -> CaO(s) + CO2(g)}. A \SI{152}{\gram} sample of \ce{CaCO3}
fully decomposes.

Find the volume of \ce{CO2} gas produced, measured at STP.
\workspace[4.5cm]
\keyonly{\begin{apcanswer}$M(\ce{CaCO3}) = 40.08+12.01+3(16.00)
= \SI{100.09}{\gram\per\mole}$

\[n(\ce{CaCO3}) = \frac{152}{100.09} = 1.519\ \text{mol}\]

1:1 ratio with \ce{CO2}, so:

\[V(\ce{CO2}) = 1.519 \times 22.42 = \mathbf{34.0}\ \text{L}\]\end{apcanswer}}

\begin{rubric}
  \rpoint{1}{molar mass of \ce{CaCO3}}
  \rpoint{1}{moles of \ce{CaCO3}}
  \rpoint{1}{uses the 1:1 stoichiometric ratio}
  \rpoint{1}{\SI{34.0}{\liter} via the STP molar volume}
\end{rubric}

% ---------------------------------------------------------------------------
\frq{long}{5}{Zumdahl \S5.4 \textbullet\ gas stoichiometry, limiting
  reactant \emph{(after Zumdahl Example 5.13)}}

\SI{2.80}{\liter} \ce{CH4} gas at \SI{25}{\celsius}, \SI{1.65}{\atm} is
mixed with \SI{35.0}{\liter} \ce{O2} gas at \SI{31}{\celsius},
\SI{1.25}{\atm} and ignited: \ce{CH4(g) + 2O2(g) -> CO2(g) + 2H2O(g)}.

Find the volume of \ce{CO2} produced, measured at \SI{2.50}{\atm} and
\SI{125}{\celsius}.
\workspace[6.5cm]
\keyonly{\begin{apcanswer}\textbf{Step 1 --- moles of each reactant.}

\[n(\ce{CH4}) = \frac{(1.65)(2.80)}{(0.08206)(298)} = 0.1889\ \text{mol}\]

\[n(\ce{O2}) = \frac{(1.25)(35.0)}{(0.08206)(304)} = 1.754\ \text{mol}\]

\textbf{Step 2 --- identify the limiting reactant.} Fully reacting
\SI{0.1889}{\mole} \ce{CH4} needs $2\times0.1889=0.3779$~mol \ce{O2}, and
\SI{1.754}{\mole} is available --- \textbf{\ce{CH4} is limiting}.

\textbf{Step 3 --- moles, then volume, of \ce{CO2}} (1:1 with \ce{CH4}):

\[n(\ce{CO2}) = 0.1889\ \text{mol}\]

\[V(\ce{CO2}) = \frac{nRT}{P} = \frac{(0.1889)(0.08206)(398)}{2.50}
= \mathbf{2.47}\ \text{L}\]\end{apcanswer}}

\begin{rubric}
  \rpoint{1}{$n(\ce{CH4})\approx\SI{0.189}{\mole}$}
  \rpoint{1}{$n(\ce{O2})\approx\SI{1.75}{\mole}$ available}
  \rpoint{1}{identifies \ce{CH4} as limiting, with the comparison shown}
  \rpoint{1}{$n(\ce{CO2})=\SI{0.189}{\mole}$ from the 1:1 ratio}
  \rpoint{1}{\SI{2.47}{\liter} at the new conditions}
\end{rubric}

% ---------------------------------------------------------------------------
\frq{short}{3}{Zumdahl \S5.4 \textbullet\ gas density and molar mass
  \emph{(after Zumdahl Example 5.14)}}

A gas sample has a density of \SI{1.95}{\gram\per\liter} at \SI{1.50}{\atm}
and \SI{27}{\celsius}.

Find its molar mass.
\workspace[3.5cm]
\keyonly{\begin{apcanswer}\[M = \frac{dRT}{P}
= \frac{(1.95)(0.08206)(300)}{1.50} = \mathbf{32.0}\ \text{g/mol}\]

(Close enough to \ce{O2}'s molar mass, \SI{32.00}{\gram\per\mole}, to be a
strong hint at the gas's identity.)\end{apcanswer}}

\begin{rubric}
  \rpoint{1}{converts to \SI{300}{\kelvin}}
  \rpoint{1}{uses $M=dRT/P$}
  \rpoint{1}{\SI{32.0}{\gram\per\mole}}
\end{rubric}

% ---------------------------------------------------------------------------
\frq{long}{5}{Zumdahl \S5.5 \textbullet\ Dalton's law from moles
  \emph{(after Zumdahl Example 5.15)}}

For a deep dive, a scuba tank is filled by pumping \SI{46}{\liter} \ce{He}
gas (at \SI{25}{\celsius}, \SI{1.0}{\atm}) and \SI{12}{\liter} \ce{O2} gas
(at \SI{25}{\celsius}, \SI{1.0}{\atm}) into a previously evacuated
\SI{5.0}{\liter} tank.

Find the partial pressure of each gas, and the total pressure, in the tank
at \SI{25}{\celsius}.
\workspace[6.5cm]
\keyonly{\begin{apcanswer}\textbf{Step 1 --- moles of each gas}, from its
own original container:

\[n(\ce{He}) = \frac{(1.0)(46)}{(0.08206)(298)} = 1.881\ \text{mol}\]

\[n(\ce{O2}) = \frac{(1.0)(12)}{(0.08206)(298)} = 0.4907\ \text{mol}\]

\textbf{Step 2 --- each gas's own partial pressure} in the shared
\SI{5.0}{\liter} tank:

\[P(\ce{He}) = \frac{nRT}{V} = \frac{(1.881)(0.08206)(298)}{5.0}
= \mathbf{9.2}\ \text{atm}\]

\[P(\ce{O2}) = \frac{(0.4907)(0.08206)(298)}{5.0} = \mathbf{2.4}\ \text{atm}\]

\textbf{Step 3 --- Dalton's law:}

\[P_{\text{total}} = P(\ce{He}) + P(\ce{O2}) = 9.2+2.4
= \mathbf{11.6}\ \text{atm}\]

(The textbook's own worked solution rounds $n(\ce{He})$ to 1.9~mol
\emph{before} recomputing pressure, which gives 9.3~atm and a total of
11.7~atm --- a good illustration of why this paper's front-matter rule is
to carry full precision and round only at the very end.)\end{apcanswer}}

\begin{rubric}
  \rpoint{1}{$n(\ce{He})\approx\SI{1.88}{\mole}$}
  \rpoint{1}{$P(\ce{He})=\SI{9.2}{\atm}$}
  \rpoint{1}{$n(\ce{O2})\approx\SI{0.491}{\mole}$}
  \rpoint{1}{$P(\ce{O2})=\SI{2.4}{\atm}$}
  \rpoint{1}{$P_{\text{total}}=\SI{11.6}{\atm}$ (9.3/2.4/11.7 accepted too,
    if a student follows the textbook's own rounded-intermediate method)}
\end{rubric}

% ---------------------------------------------------------------------------
\frq{short}{2}{Zumdahl \S5.5 \textbullet\ Dalton's law, mole fraction
  \emph{(after Zumdahl Example 5.16)}}

In a sample of air, \ce{O2}'s partial pressure is \SI{156}{\torr} while
the total pressure is \SI{743}{\torr}.

Find the mole fraction of \ce{O2}.
\workspace[2.5cm]
\keyonly{\begin{apcanswer}\[\chi(\ce{O2}) = \frac{P(\ce{O2})}
{P_{\text{total}}} = \frac{156}{743} = \mathbf{0.210}\]\end{apcanswer}}

\begin{rubric}
  \rpoint{1}{uses $\chi_A = P_A/P_{\text{total}}$}
  \rpoint{1}{0.210}
\end{rubric}

% ---------------------------------------------------------------------------
\frq{short}{2}{Zumdahl \S5.5 \textbullet\ Dalton's law, partial pressure
  from mole fraction \emph{(after Zumdahl Example 5.17)}}

Air is $0.7808$ mole fraction \ce{N2}.

Find the partial pressure of \ce{N2} when the total pressure is
\SI{760.}{\torr}.
\workspace[2.5cm]
\keyonly{\begin{apcanswer}\[P(\ce{N2}) = \chi(\ce{N2}) \times
P_{\text{total}} = 0.7808\times760 = \mathbf{593}\ \text{torr}\]
\end{apcanswer}}

\begin{rubric}
  \rpoint{1}{uses $P_A=\chi_AP_{\text{total}}$}
  \rpoint{1}{\SI{593}{\torr}}
\end{rubric}

% ---------------------------------------------------------------------------
\frq{long}{5}{Zumdahl \S5.5 \textbullet\ gas collection over water
  \emph{(after Zumdahl Example 5.18)}}

\ce{2KClO3(s) -> 2KCl(s) + 3O2(g)}. The \ce{O2} produced is collected by
water displacement at \SI{22}{\celsius}: the collected gas has a total
pressure of \SI{754}{\torr} and a volume of \SI{0.650}{\liter}. Water's
vapor pressure at \SI{22}{\celsius} is \SI{21}{\torr}.

\begin{parts}
  \item Find the partial pressure of \ce{O2} in the collected gas.
        \workspace[3.5cm]
        \keyonly{\begin{apcanswer}$754-21=\mathbf{733}$~torr $=\SI{0.9645}{\atm}$\end{apcanswer}}
  \item Find the mass of \ce{KClO3} that originally decomposed.
        \workspace[5cm]
        \keyonly{\begin{apcanswer}\[n(\ce{O2}) = \frac{PV}{RT}
        = \frac{(0.9645)(0.650)}{(0.08206)(295)} = 0.02590\ \text{mol}\]

        2:3 ratio with \ce{KClO3}:

        \[n(\ce{KClO3}) = 0.02590\times\frac{2}{3} = 0.01726\ \text{mol}\]

        $M(\ce{KClO3}) = 39.10+35.45+3(16.00)
        = \SI{122.55}{\gram\per\mole}$

        \[\text{mass} = 0.01726\times122.55 = \mathbf{2.12}\ \text{g}\]
        \end{apcanswer}}
\end{parts}

\begin{rubric}
  \rpoint{1}{(a) subtracts water's vapor pressure to reach \SI{733}{\torr}}
  \rpoint{1}{(b) converts pressure to atm}
  \rpoint{1}{(b) finds $n(\ce{O2})\approx\SI{0.0259}{\mole}$}
  \rpoint{1}{(b) applies the 2:3 stoichiometric ratio}
  \rpoint{1}{(b) \SI{2.12}{\gram} using the correct molar mass of
    \ce{KClO3}}
\end{rubric}

% ---------------------------------------------------------------------------
\frq{short}{3}{Zumdahl \S5.7 \textbullet\ root-mean-square velocity
  \textbf{[beyond CED]} \emph{(after Zumdahl Example 5.19)}}

Calculate the root-mean-square velocity of \ce{He} gas atoms at
\SI{25}{\celsius}.
\workspace[4cm]
\keyonly{\begin{apcanswer}Use SI units: $M(\ce{He})=\SI{4.003e-3}
{\kilo\gram\per\mole}$, $R=\SI{8.3145}{\joule\per\mole\per\kelvin}$,
$T=\SI{298}{\kelvin}$.

\[u_{\text{rms}} = \sqrt{\frac{3RT}{M}}
= \sqrt{\frac{3(8.3145)(298)}{4.003\times10^{-3}}}
= \mathbf{1.36\times10^{3}}\ \text{m/s}\]\end{apcanswer}}

\begin{rubric}
  \rpoint{1}{correct formula $u_{\text{rms}}=\sqrt{3RT/M}$ with $R$ in
    \si{\joule\per\mole\per\kelvin} and $M$ in \si{\kilo\gram\per\mole}}
  \rpoint{1}{correct substitution}
  \rpoint{1}{\SI{1.36e3}{\meter\per\second}}
\end{rubric}

% ---------------------------------------------------------------------------
\frq{short}{3}{Zumdahl \S5.7 \textbullet\ Graham's law of effusion
  \textbf{[beyond CED]} \emph{(after Zumdahl Example 5.20)}}

Calculate the ratio of effusion rates of \ce{H2} gas to \ce{UF6} gas,
under identical conditions.
\workspace[4cm]
\keyonly{\begin{apcanswer}\[\frac{\text{rate}(\ce{H2})}
{\text{rate}(\ce{UF6})} = \sqrt{\frac{M(\ce{UF6})}{M(\ce{H2})}}
= \sqrt{\frac{352.03}{2.016}} = \mathbf{13.2}\]

Hydrogen effuses about 13 times faster than uranium hexafluoride. Real
gaseous-diffusion uranium enrichment uses this same Graham's-law idea,
but on a far more subtle scale: it separates $^{235}\ce{UF6}$ from
$^{238}\ce{UF6}$ (two isotopologues of the \emph{same} compound), where
the mass difference is tiny and the per-stage separation factor is only
about 1.004 --- thousands of stages are needed. \ce{H2} vs.\ \ce{UF6}
illustrates the physics, not the actual industrial
process.\end{apcanswer}}

\begin{rubric}
  \rpoint{1}{correct Graham's law setup}
  \rpoint{1}{correct molar masses substituted}
  \rpoint{1}{ratio $\approx13.2$}
\end{rubric}

% ============================================================================
% SESSION 2
% ============================================================================
\clearpage
\examsection{C}{Gas Laws}{24 questions \textbullet\ 73 points}{\calculatorok}

% ---------------------------------------------------------------------------
\frq{short}{3}{Zumdahl \S5.1 \textbullet\ pressure conversion
  \emph{(after Zumdahl Q42)}}

A gas cylinder's gauge reads 2200~psi ($\SI{1}{\atm}=14.7$~psi).

Convert this pressure to (a) atm, (b) MPa, (c) torr.
\workspace[4cm]
\keyonly{\begin{apcanswer}$2200/14.7 = \mathbf{150}$~atm

$150\times101325 = 1.5\times10^7$~Pa $= \mathbf{15}$~MPa

$150\times760 = \mathbf{1.1\times10^5}$~torr\end{apcanswer}}

\begin{rubric}
  \rpoint{1}{\SI{150}{\atm}}
  \rpoint{1}{\SI{15}{\mega\pascal}}
  \rpoint{1}{\SI{1.1e5}{\torr}}
\end{rubric}

% ---------------------------------------------------------------------------
\frq{long}{4}{Zumdahl \S5.1 \textbullet\ manometers and fluid choice
  \emph{(after Zumdahl Q46)}}

A gas sample is connected to an open-tube mercury manometer; atmospheric
pressure is \SI{760.}{\torr}. In two separate trials, the gas pressure
exceeds atmospheric pressure, and the mercury column shows height
differences of \SI{118}{\milli\meter} and \SI{215}{\milli\meter}.

\begin{parts}
  \item Find the flask pressure in each trial, in torr, atm, and Pa.
        \workspace[5cm]
        \keyonly{\begin{apcanswer}Flask pressure $=P_{\text{atm}}+\Delta h$:

        Trial 1: $760+118=\mathbf{878}$~torr
        $=\SI{1.16}{\atm}=\SI{1.17e5}{\pascal}$

        Trial 2: $760+215=\mathbf{975}$~torr
        $=\SI{1.28}{\atm}=\SI{1.30e5}{\pascal}$\end{apcanswer}}
  \item If the manometer were instead filled with nonvolatile silicone oil
        (density \SI{1.30}{\gram\per\centi\meter\cubed}) rather than
        mercury (density \SI{13.6}{\gram\per\centi\meter\cubed}), would
        the observed column-height difference for the \emph{same} gas
        pressure be larger or smaller? Explain.
        \workspace[4.5cm]
        \keyonly{\begin{apcanswer}Larger. Column height is inversely proportional
        to fluid density at a fixed pressure difference --- mercury is
        about 10 times denser than this oil, so the same pressure
        difference would push an oil column roughly 10 times higher. A
        taller column is easier to read precisely, which is exactly why a
        less-dense fluid is chosen for measuring small pressure
        differences.\end{apcanswer}}
\end{parts}

\begin{rubric}
  \rpoint{1}{Trial 1: \SI{878}{\torr} (or equivalent atm/Pa)}
  \rpoint{1}{Trial 2: \SI{975}{\torr} (or equivalent atm/Pa)}
  \rpoint{2}{explains that a less dense fluid gives a taller, more
    readable column for the same pressure, via height $\propto1/$density}
\end{rubric}

% ---------------------------------------------------------------------------
\frq{short}{2}{Zumdahl \S5.2 \textbullet\ Boyle's law
  \emph{(after Zumdahl Q47)}}

A \SI{400.}{\milli\liter} can of compressed gas at \SI{5.20}{\atm} is
sprayed completely into a large, previously empty plastic bag, which
inflates to \SI{2.14}{\liter} (constant temperature).

Find the gas pressure in the bag.
\workspace[3cm]
\keyonly{\begin{apcanswer}\[P_2 = \frac{P_1V_1}{V_2}
= \frac{(5.20)(0.400)}{2.14} = \mathbf{0.972}\ \text{atm}\]\end{apcanswer}}

\begin{rubric}
  \rpoint{1}{sets up Boyle's law with matched volume units}
  \rpoint{1}{\SI{0.972}{\atm}}
\end{rubric}

% ---------------------------------------------------------------------------
\frq{short}{2}{Zumdahl \S5.3 \textbullet\ Charles's law
  \emph{(after Zumdahl Q48)}}

A balloon is filled to $7.00\times10^2$~mL at \SI{20.0}{\celsius}, then
cooled at constant pressure to \SI{1.00e2}{\kelvin}.

Find the final volume.
\workspace[3cm]
\keyonly{\begin{apcanswer}$T_1=\SI{293}{\kelvin}$

\[V_2 = 700.\times\frac{100.}{293} = \mathbf{239}\ \text{mL}\]
\end{apcanswer}}

\begin{rubric}
  \rpoint{1}{converts \SI{20.0}{\celsius} to \SI{293}{\kelvin}}
  \rpoint{1}{\SI{239}{\milli\liter}}
\end{rubric}

% ---------------------------------------------------------------------------
\frq{short}{3}{Zumdahl \S5.3 \textbullet\ Avogadro's law from molar volume
  \emph{(after Zumdahl Q49)}}

\SI{27.1}{\gram} \ce{Ar(g)} occupies \SI{4.21}{\liter}.

What volume would \SI{1.29}{\mole} \ce{Ne(g)} occupy at the same
temperature and pressure?
\workspace[4cm]
\keyonly{\begin{apcanswer}\[n(\ce{Ar}) = \frac{27.1}{39.95}
= 0.6783\ \text{mol} \ \Longrightarrow\
V_m = \frac{4.21}{0.6783} = 6.206\ \text{L/mol}\]

\[V(\ce{Ne}) = 1.29\times6.206 = \mathbf{8.01}\ \text{L}\]\end{apcanswer}}

\begin{rubric}
  \rpoint{1}{finds moles of \ce{Ar}}
  \rpoint{1}{finds the shared molar volume at these conditions}
  \rpoint{1}{\SI{8.01}{\liter}}
\end{rubric}

% ---------------------------------------------------------------------------
\frq{short}{2}{Zumdahl \S5.3 \textbullet\ stoichiometric volume ratios
  \emph{(after Zumdahl Q50)}}

\ce{2NO2(g) -> N2O4(g)}. If \SI{25.0}{\milli\liter} \ce{NO2} gas fully
converts to \ce{N2O4} under the same conditions, what volume does the
\ce{N2O4} occupy?
\workspace[2.5cm]
\keyonly{\begin{apcanswer}Gas volumes track moles at fixed $T,P$, and the
2:1 ratio halves the volume:

\[V(\ce{N2O4}) = 25.0\times\frac{1}{2} = \mathbf{12.5}\ \text{mL}\]
\end{apcanswer}}

\begin{rubric}
  \rpoint{1}{recognizes $V\propto n$ applies here}
  \rpoint{1}{\SI{12.5}{\milli\liter}}
\end{rubric}

% ---------------------------------------------------------------------------
\frq{long}{4}{Zumdahl \S5.4 \textbullet\ complete the table I
  \emph{(after Zumdahl Q51)}}

Complete the table below for an ideal gas ($R=\SI{0.08206}{\liter\atm
\per\mole\per\kelvin}$):

\begin{center}
\begin{tabular}{ccccc}
\toprule
 & $P$ (atm) & $V$ (L) & $n$ (mol) & $T$ \\
\midrule
a. & 5.00 & --- & 2.00 & \SI{155}{\celsius} \\
b. & 0.300 & 2.00 & --- & \SI{155}{\kelvin} \\
c. & 4.47 & 25.0 & 2.01 & --- \\
d. & --- & 2.25 & 10.5 & \SI{75}{\celsius} \\
\bottomrule
\end{tabular}
\end{center}
\workspace[5cm]
\keyonly{\begin{apcanswer}\textbf{a.} $T=\SI{428}{\kelvin}$;
$V=nRT/P=(2.00)(0.08206)(428)/5.00=\mathbf{14.1}$~L

\textbf{b.} $n=PV/RT=(0.300)(2.00)/[(0.08206)(155)]=\mathbf{0.0472}$~mol

\textbf{c.} $T=PV/nR=(4.47)(25.0)/[(2.01)(0.08206)]=\mathbf{678}$~K

\textbf{d.} $T=\SI{348}{\kelvin}$;
$P=nRT/V=(10.5)(0.08206)(348)/2.25=\mathbf{133}$~atm\end{apcanswer}}

\begin{rubric}
  \rpoint{1}{a. \SI{14.1}{\liter}}
  \rpoint{1}{b. \SI{0.0472}{\mole}}
  \rpoint{1}{c. \SI{678}{\kelvin}}
  \rpoint{1}{d. \SI{133}{\atm}}
\end{rubric}

% ---------------------------------------------------------------------------
\frq{long}{4}{Zumdahl \S5.4 \textbullet\ complete the table II
  \emph{(after Zumdahl Q52)}}

Complete the table below:

\begin{center}
\begin{tabular}{ccccc}
\toprule
 & $P$ & $V$ & $n$ (mol) & $T$ \\
\midrule
a. & $7.74\times10^3$~Pa & 12.2 mL & --- & \SI{25}{\celsius} \\
b. & --- & 43.0 mL & 0.421 & \SI{223}{\kelvin} \\
c. & 455 torr & --- & $4.4\times10^{-2}$ & \SI{331}{\celsius} \\
d. & 745 mmHg & 11.2 L & 0.401 & --- \\
\bottomrule
\end{tabular}
\end{center}
\workspace[5cm]
\keyonly{\begin{apcanswer}\textbf{a.} $P=0.07639$~atm, $V=0.0122$~L,
$T=298$~K: $n=PV/RT=\mathbf{3.81\times10^{-5}}$~mol

\textbf{b.} $P=nRT/V=(0.421)(0.08206)(223)/0.0430
=\mathbf{179}$~atm

\textbf{c.} $P=0.5987$~atm, $T=604$~K:
$V=nRT/P=\mathbf{3.6}$~L

\textbf{d.} $P=0.9803$~atm: $T=PV/nR=(0.9803)(11.2)/[(0.401)(0.08206)]
=\mathbf{334}$~K\end{apcanswer}}

\begin{rubric}
  \rpoint{1}{a. $n\approx3.81\times10^{-5}$~mol}
  \rpoint{1}{b. $P\approx\SI{179}{\atm}$}
  \rpoint{1}{c. $V\approx\SI{3.6}{\liter}$}
  \rpoint{1}{d. $T\approx\SI{334}{\kelvin}$}
\end{rubric}

% ---------------------------------------------------------------------------
\frq{long}{4}{Zumdahl \S5.4 \textbullet\ moles of air in the lungs
  \emph{(after Zumdahl Q54)}}

Average human lung capacity is \SI{6.0}{\liter}. Find the moles of air in
the lungs:

\begin{parts}
  \item At sea level: $T=\SI{298}{\kelvin}$, $P=\SI{1.00}{\atm}$.
        \workspace[2.5cm]
        \keyonly{\begin{apcanswer}$n=PV/RT=\mathbf{0.25}$~mol\end{apcanswer}}
  \item \SI{10.}{\meter} underwater: $T=\SI{298}{\kelvin}$,
        $P=\SI{1.97}{\atm}$.
        \workspace[2.5cm]
        \keyonly{\begin{apcanswer}$n=\mathbf{0.48}$~mol\end{apcanswer}}
  \item Summit of Mount Everest: $T=\SI{200.}{\kelvin}$,
        $P=\SI{0.296}{\atm}$.
        \workspace[2.5cm]
        \keyonly{\begin{apcanswer}$n=\mathbf{0.11}$~mol\end{apcanswer}}
\end{parts}

\begin{rubric}
  \rpoint{1}{(a) \SI{0.25}{\mole}}
  \rpoint{1}{(b) \SI{0.48}{\mole}}
  \rpoint{2}{(c) \SI{0.11}{\mole} --- both the low $T$ and low $P$ used
    correctly}
\end{rubric}

% ---------------------------------------------------------------------------
\frq{short}{2}{Zumdahl \S5.4 \textbullet\ ideal gas law, mass given
  \emph{(after Zumdahl Q56)}}

A \SI{5.0}{\liter} flask holds \SI{0.60}{\gram} \ce{O2} at
\SI{22}{\celsius}.

Find the pressure, in atm.
\workspace[3cm]
\keyonly{\begin{apcanswer}$n=0.60/32.00=0.01875$~mol, $T=295$~K

\[P = \frac{nRT}{V} = \frac{(0.01875)(0.08206)(295)}{5.0}
= \mathbf{0.091}\ \text{atm}\]\end{apcanswer}}

\begin{rubric}
  \rpoint{1}{converts mass to moles}
  \rpoint{1}{\SI{0.091}{\atm}}
\end{rubric}

% ---------------------------------------------------------------------------
\frq{short}{3}{Zumdahl \S5.4 \textbullet\ ideal gas law, two unknowns
  \emph{(after Zumdahl Q57)}}

A \SI{2.50}{\liter} container holds \SI{175}{\gram} argon.

\begin{parts}
  \item At \SI{10.0}{\atm}, find the temperature.
        \workspace[2.5cm]
        \keyonly{\begin{apcanswer}$n=4.380$~mol; $T=PV/nR=\mathbf{69.5}$~K\end{apcanswer}}
  \item At \SI{225}{\kelvin}, find the pressure.
        \workspace[2.5cm]
        \keyonly{\begin{apcanswer}$P=nRT/V=\mathbf{32.4}$~atm\end{apcanswer}}
\end{parts}

\begin{rubric}
  \rpoint{1}{moles of argon, \SI{4.38}{\mole}}
  \rpoint{1}{(a) \SI{69.5}{\kelvin} (note: this is below argon's triple
    point of \SI{83.8}{\kelvin}, so real argon would already be solid ---
    this is the idealized $PV=nRT$ result)}
  \rpoint{1}{(b) \SI{32.4}{\atm}}
\end{rubric}

% ---------------------------------------------------------------------------
\frq{short}{3}{Zumdahl \S5.4 \textbullet\ liquid-to-gas volume
  \emph{(after Zumdahl Q58)}}

A person swallows a \SI{0.050}{\milli\liter} drop of liquid \ce{O2}
(density \SI{1.149}{\gram\per\milli\liter}).

What volume of gas results in the stomach, at \SI{37}{\celsius} and
\SI{1.0}{\atm}?
\workspace[4cm]
\keyonly{\begin{apcanswer}mass $=0.050\times1.149=0.05745$~g,
$n=0.05745/32.00=1.795\times10^{-3}$~mol, $T=310$~K

\[V = \frac{nRT}{P} = \mathbf{0.046}\ \text{L} = 46\ \text{mL}\]
\end{apcanswer}}

\begin{rubric}
  \rpoint{1}{converts the liquid volume to a mass, then to moles}
  \rpoint{1}{correct kelvin temperature}
  \rpoint{1}{\SI{46}{\milli\liter}}
\end{rubric}

% ---------------------------------------------------------------------------
\frq{short}{3}{Zumdahl \S5.4 \textbullet\ counting gas particles
  \emph{(after Zumdahl Q59)}}

A breath draws in \SI{450}{\milli\liter} of air at \SI{745}{\torr},
\SI{22}{\celsius}.

How many air particles are in the breath?
\workspace[4cm]
\keyonly{\begin{apcanswer}$P=745/760=0.9803$~atm, $T=295$~K

\[n = \frac{PV}{RT} = \frac{(0.9803)(0.450)}{(0.08206)(295)}
= 0.01822\ \text{mol}\]

\[\text{particles} = 0.01822\times6.022\times10^{23}
= \mathbf{1.10\times10^{22}}\]\end{apcanswer}}

\begin{rubric}
  \rpoint{1}{converts torr to atm}
  \rpoint{1}{finds $n\approx0.0182$~mol}
  \rpoint{1}{$1.10\times10^{22}$ particles}
\end{rubric}

% ---------------------------------------------------------------------------
\frq{short}{3}{Zumdahl \S5.4 \textbullet\ transferring gas between
  containers \emph{(after Zumdahl Q60)}}

\ce{N2O} fills a cylinder at \SI{10.5}{\atm}, \SI{5.00}{\liter},
\SI{298}{\kelvin}.

\begin{parts}
  \item Find the moles of \ce{N2O} present.
        \workspace[2.5cm]
        \keyonly{\begin{apcanswer}$n=\mathbf{2.15}$~mol\end{apcanswer}}
  \item That gas is emptied into a large balloon at \SI{745}{\torr}. Find
        the balloon's volume at \SI{298}{\kelvin}.
        \workspace[3.5cm]
        \keyonly{\begin{apcanswer}$P=0.9803$~atm;
        $V=nRT/P=\mathbf{53.6}$~L\end{apcanswer}}
\end{parts}

\begin{rubric}
  \rpoint{1}{(a) \SI{2.15}{\mole}}
  \rpoint{1}{(b) converts \SI{745}{\torr} to atm}
  \rpoint{1}{(b) \SI{53.6}{\liter}}
\end{rubric}

% ---------------------------------------------------------------------------
\frq{long}{4}{Zumdahl \S5.4 \textbullet\ gas added to a rigid container
  \emph{(after Zumdahl Q61)}}

A rigid container holds \SI{1.50}{\mole} gas at \SI{25}{\celsius},
exerting \SI{400.}{\torr}. More gas is added and the temperature raised to
\SI{50.}{\celsius}; the pressure rises to \SI{800.}{\torr}.

How many moles were added?
\workspace[5.5cm]
\keyonly{\begin{apcanswer}\textbf{Step 1 --- find the fixed container
volume} from the initial state:

\[V = \frac{n_1RT_1}{P_1} = \frac{(1.50)(0.08206)(298)}{400/760}
= 69.7\ \text{L}\]

\textbf{Step 2 --- find the final moles}, same $V$, new $P,T$:

\[n_2 = \frac{P_2V}{RT_2} = \frac{(800/760)(69.7)}{(0.08206)(323)}
= 2.77\ \text{mol}\]

\textbf{Step 3 --- subtract:}

\[\Delta n = 2.77-1.50 = \mathbf{1.27}\ \text{mol added}\]\end{apcanswer}}

\begin{rubric}
  \rpoint{1}{finds the container's fixed volume from the initial state}
  \rpoint{1}{converts torr to atm for both states}
  \rpoint{1}{finds $n_2\approx\SI{2.77}{\mole}$}
  \rpoint{1}{\SI{1.27}{\mole} added}
\end{rubric}

% ---------------------------------------------------------------------------
\frq{short}{3}{Zumdahl \S5.4 \textbullet\ combined gas law, percent change
  \emph{(after Zumdahl Q62)}}

A bicycle tire is inflated to 75~psi at \SI{19}{\celsius}. Riding
on hot asphalt raises the tire to \SI{58}{\celsius}, and its volume grows
4.0\%.

Find the new pressure.
\workspace[4cm]
\keyonly{\begin{apcanswer}$T_1=292$~K, $T_2=331$~K,
$V_2/V_1=1.040$

\[P_2 = P_1\times\frac{V_1}{V_2}\times\frac{T_2}{T_1}
= 75\times\frac{1}{1.040}\times\frac{331}{292}
= \mathbf{82}\ \text{psi}\]\end{apcanswer}}

\begin{rubric}
  \rpoint{1}{converts both temperatures to kelvin}
  \rpoint{1}{accounts for the 4.0\% volume increase}
  \rpoint{1}{82 psi}
\end{rubric}

% ---------------------------------------------------------------------------
\frq{short}{3}{Zumdahl \S5.4 \textbullet\ comparing two containers
  \emph{(after Zumdahl Q63)}}

Container A: \ce{SO2(g)}, pressure $P_A$, \SI{1.0}{\mole},
\SI{1.0}{\liter}, \SI{7}{\celsius}.

Container B: an unknown gas, pressure $P_B$, \SI{2.0}{\mole},
\SI{2.0}{\liter}, \SI{-87}{\celsius}.

How does $P_B$ relate to $P_A$?
\workspace[4.5cm]
\keyonly{\begin{apcanswer}\[P_A = \frac{n_ART_A}{V_A}
= \frac{(1.0)(0.08206)(280)}{1.0} = 23.0\ \text{atm}\]

\[P_B = \frac{n_BRT_B}{V_B}
= \frac{(2.0)(0.08206)(186)}{2.0} = 15.3\ \text{atm}\]

Since $n/V=1.0$~mol/L for both containers, this simplifies to
$P_B/P_A = T_B/T_A = 186/280$:

\[\mathbf{P_B \approx 0.66\,P_A}\]\end{apcanswer}}

\begin{rubric}
  \rpoint{1}{$P_A\approx\SI{23.0}{\atm}$}
  \rpoint{1}{$P_B\approx\SI{15.3}{\atm}$}
  \rpoint{1}{states the relationship, $P_B\approx0.66P_A$}
\end{rubric}

% ---------------------------------------------------------------------------
\frq{short}{2}{Zumdahl \S5.4 \textbullet\ conceptual, simultaneous changes
  \emph{(after Zumdahl Q64)}}

What happens to an ideal gas's volume if pressure doubles while the
absolute temperature is simultaneously halved?
\workspace[3cm]
\keyonly{\begin{apcanswer}Since $V\propto T/P$:

\[\frac{V_2}{V_1} = \frac{T_2/T_1}{P_2/P_1} = \frac{0.5}{2}
= \mathbf{0.25}\]

The volume drops to one quarter of its original value.\end{apcanswer}}

\begin{rubric}
  \rpoint{1}{sets up $V\propto T/P$}
  \rpoint{1}{correctly concludes $V_2=\tfrac14 V_1$}
\end{rubric}

% ---------------------------------------------------------------------------
\frq{long}{3}{Zumdahl \S5.4 \textbullet\ Gay-Lussac's law
  \emph{(after Zumdahl Q65)}}

A rigid container holds ideal gas at \SI{11.0}{\atm}, \SI{0}{\celsius}.

\begin{parts}
  \item Find the pressure if heated to \SI{45}{\celsius}.
        \workspace[2.5cm]
        \keyonly{\begin{apcanswer}$\mathbf{12.8}$~atm\end{apcanswer}}
  \item Find the temperature for a pressure of \SI{6.50}{\atm}.
        \workspace[2.5cm]
        \keyonly{\begin{apcanswer}$\mathbf{161}$~K\end{apcanswer}}
  \item Find the temperature for a pressure of \SI{25.0}{\atm}.
        \workspace[2.5cm]
        \keyonly{\begin{apcanswer}$\mathbf{620}$~K\end{apcanswer}}
\end{parts}

\begin{rubric}
  \rpoint{1}{(a) \SI{12.8}{\atm}}
  \rpoint{1}{(b) \SI{161}{\kelvin}}
  \rpoint{1}{(c) \SI{620}{\kelvin}}
\end{rubric}

% ---------------------------------------------------------------------------
\frq{short}{3}{Zumdahl \S5.4 \textbullet\ Charles's law in three dimensions
  \emph{(after Zumdahl Q66)}}

An ideal gas at \SI{7}{\celsius} fills a spherical, flexible container of
radius \SI{1.00}{\centi\meter}. Heated at constant pressure to
\SI{88}{\celsius}.

Find the new radius. ($V_{\text{sphere}}=\tfrac43\pi r^3$.)
\workspace[4cm]
\keyonly{\begin{apcanswer}$T_1=280$~K, $T_2=361$~K; $V\propto T$, so
$V_2/V_1=T_2/T_1=1.289$.

Since $V\propto r^3$:

\[r_2 = r_1\times(1.289)^{1/3} = \mathbf{1.09}\ \text{cm}\]
\end{apcanswer}}

\begin{rubric}
  \rpoint{1}{converts both temperatures to kelvin}
  \rpoint{1}{finds the volume ratio $V_2/V_1=1.29$}
  \rpoint{1}{takes the cube root correctly, \SI{1.09}{\centi\meter}}
\end{rubric}

% ---------------------------------------------------------------------------
\frq{short}{3}{Zumdahl \S5.4 \textbullet\ extreme compression
  \emph{(after Zumdahl Q67)}}

An ideal gas in a $5.00\times10^2$~mL cylinder at \SI{30.}{\celsius} and
\SI{710.}{\torr} is compressed to \SI{25}{\milli\liter} while heated to
\SI{820.}{\celsius}.

Find the new pressure.
\workspace[4cm]
\keyonly{\begin{apcanswer}$T_1=303$~K, $T_2=1093$~K

\[P_2 = P_1\times\frac{V_1}{V_2}\times\frac{T_2}{T_1}
= 710.\times\frac{500.}{25}\times\frac{1093}{303}
= \mathbf{5.12\times10^4}\ \text{torr}\]

($\approx\SI{67.4}{\atm}$ --- a reminder that the combined gas law does
not know or care that no real cylinder survives this.)\end{apcanswer}}

\begin{rubric}
  \rpoint{1}{converts both temperatures to kelvin}
  \rpoint{1}{sets up the combined gas law correctly}
  \rpoint{1}{$5.12\times10^4$~torr}
\end{rubric}

% ---------------------------------------------------------------------------
\frq{long}{4}{Zumdahl \S5.4 \textbullet\ mass remaining in a cylinder
  \emph{(after Zumdahl Q68)}}

A compressed-gas cylinder holds $1.00\times10^3$~g argon at
2050.~psi, \SI{18}{\celsius}.

How much gas (mass) remains if the pressure drops to 650.~psi at
\SI{26}{\celsius}?
\workspace[6cm]
\keyonly{\begin{apcanswer}\textbf{Step 1 --- find the cylinder's fixed
volume} from the initial state ($n_1=1000/39.95=25.03$~mol,
$P_1=2050/14.7=139.5$~atm, $T_1=291$~K):

\[V = \frac{n_1RT_1}{P_1} = \frac{(25.03)(0.08206)(291)}{139.5}
= 4.29\ \text{L}\]

\textbf{Step 2 --- find the remaining moles} at the new state
($P_2=650/14.7=44.2$~atm, $T_2=299$~K):

\[n_2 = \frac{P_2V}{RT_2} = \frac{(44.2)(4.29)}{(0.08206)(299)}
= 7.72\ \text{mol}\]

\textbf{Step 3 --- convert back to mass:}

\[\text{mass} = 7.72\times39.95 = \mathbf{309}\ \text{g}\]\end{apcanswer}}

\begin{rubric}
  \rpoint{1}{finds the cylinder's fixed volume from the initial state}
  \rpoint{1}{converts psi to atm for both states}
  \rpoint{1}{finds $n_2\approx\SI{7.72}{\mole}$}
  \rpoint{1}{\SI{309}{\gram} remaining}
\end{rubric}

% ---------------------------------------------------------------------------
\frq{short}{3}{Zumdahl \S5.4 \textbullet\ balloon at altitude
  \emph{(after Zumdahl Q69)}}

A sealed balloon holds \SI{1.00}{\liter} \ce{He} at \SI{23}{\celsius},
\SI{1.00}{\atm}. It rises to where the pressure is \SI{220.}{\torr} and
the temperature is \SI{-31}{\celsius}.

Find the new volume, and the volume change from the initial \SI{1.00}
{\liter}.
\workspace[4.5cm]
\keyonly{\begin{apcanswer}$T_1=296$~K, $T_2=242$~K,
$P_2=220/760=0.2895$~atm

\[V_2 = V_1\times\frac{P_1}{P_2}\times\frac{T_2}{T_1}
= 1.00\times\frac{1.00}{0.2895}\times\frac{242}{296}
= \mathbf{2.82}\ \text{L}\]

\[\Delta V = 2.82-1.00 = \mathbf{+1.82}\ \text{L increase}\]
\end{apcanswer}}

\begin{rubric}
  \rpoint{1}{converts both temperatures to kelvin and \SI{220.}{\torr} to
    atm}
  \rpoint{1}{\SI{2.82}{\liter}}
  \rpoint{1}{correctly reports the \SI{1.82}{\liter} increase}
\end{rubric}

% ---------------------------------------------------------------------------
\frq{short}{3}{Zumdahl \S5.4 \textbullet\ open-system mole ratio
  \emph{(after Zumdahl Q70)}}

A hot-air balloon holds $4.00\times10^3$~m$^3$ air at \SI{745}{\torr},
\SI{21}{\celsius}. Heated to \SI{62}{\celsius}, it expands to
$4.20\times10^3$~m$^3$ (pressure unchanged --- the balloon is open to the
atmosphere).

Find the ratio of final to initial moles of air.
\workspace[4cm]
\keyonly{\begin{apcanswer}At constant $P$, $n\propto V/T$:

\[\frac{n_2}{n_1} = \frac{V_2/T_2}{V_1/T_1}
= \frac{4.20\times10^3/335}{4.00\times10^3/294}
= \mathbf{0.922}\]

The balloon holds about 7.8\% \emph{fewer} moles of air after heating ---
this is exactly why hot-air balloons rise: heated air is driven out
through the open bottom, so the balloon's average density drops even
though its volume grows.\end{apcanswer}}

\begin{rubric}
  \rpoint{1}{converts both temperatures to kelvin}
  \rpoint{1}{sets up $n\propto V/T$ at constant $P$}
  \rpoint{1}{$n_2/n_1\approx0.922$, with the physical interpretation}
\end{rubric}

% ============================================================================
% SESSION 3
% ============================================================================
\clearpage
\examsection{D}{Gas Density, Molar Mass, and Reaction Stoichiometry}
  {16 questions \textbullet\ 54 points}{\calculatorok}

% ---------------------------------------------------------------------------
\frq{short}{3}{Zumdahl \S5.4 \textbullet\ gas-to-mass stoichiometry
  \emph{(after Zumdahl Q71)}}

\ce{4Al(s) + 3O2(g) -> 2Al2O3(s)}. \SI{2.00}{\liter} pure \ce{O2} at STP
fully reacts with a sample of aluminum.

Find the mass of aluminum consumed.
\workspace[4cm]
\keyonly{\begin{apcanswer}\[n(\ce{O2}) = \frac{2.00}{22.42}
= 0.08921\ \text{mol}\]

4:3 ratio with \ce{Al}:

\[n(\ce{Al}) = 0.08921\times\frac43 = 0.1190\ \text{mol}\]

\[\text{mass} = 0.1190\times26.98 = \mathbf{3.21}\ \text{g}\]\end{apcanswer}}

\begin{rubric}
  \rpoint{1}{moles of \ce{O2} via STP molar volume}
  \rpoint{1}{applies the 4:3 ratio}
  \rpoint{1}{\SI{3.21}{\gram}}
\end{rubric}

% ---------------------------------------------------------------------------
\frq{short}{2}{Zumdahl \S5.4 \textbullet\ sublimation to gas volume
  \emph{(after Zumdahl Q72)}}

A student adds \SI{4.00}{\gram} dry ice (solid \ce{CO2}) to an empty
balloon.

Find the balloon's volume at STP once all the \ce{CO2} sublimes.
\workspace[3cm]
\keyonly{\begin{apcanswer}\[n = \frac{4.00}{44.01} = 0.0909\ \text{mol}
\ \Longrightarrow\ V = 0.0909\times22.42 = \mathbf{2.04}\ \text{L}\]
\end{apcanswer}}

\begin{rubric}
  \rpoint{1}{moles of \ce{CO2}}
  \rpoint{1}{\SI{2.04}{\liter}}
\end{rubric}

% ---------------------------------------------------------------------------
\frq{short}{3}{Zumdahl \S5.4 \textbullet\ airbag stoichiometry
  \emph{(after Zumdahl Q73)}}

Airbags deploy when sodium azide decomposes explosively:
\ce{2NaN3(s) -> 2Na(s) + 3N2(g)}.

What mass of \ce{NaN3(s)} must react to inflate a bag to \SI{70.0}{\liter}
\ce{N2} at STP?
\workspace[4cm]
\keyonly{\begin{apcanswer}$M(\ce{NaN3}) = 22.99+3(14.01)
= \SI{65.02}{\gram\per\mole}$

\[n(\ce{N2}) = \frac{70.0}{22.42} = 3.123\ \text{mol}\]

2:3 ratio with \ce{NaN3}:

\[n(\ce{NaN3}) = 3.123\times\frac23 = 2.082\ \text{mol}\]

\[\text{mass} = 2.082\times65.02 = \mathbf{135}\ \text{g}\]\end{apcanswer}}

\begin{rubric}
  \rpoint{1}{molar mass of \ce{NaN3}}
  \rpoint{1}{moles of \ce{N2}, then applies the 2:3 ratio}
  \rpoint{1}{\SI{135}{\gram}}
\end{rubric}

% ---------------------------------------------------------------------------
\frq{short}{3}{Zumdahl \S5.4 \textbullet\ decomposition of a solution
  \emph{(after Zumdahl Q74)}}

\ce{2H2O2(aq) -> 2H2O(l) + O2(g)}. What volume of pure \ce{O2(g)},
collected at \SI{27}{\celsius} and \SI{746}{\torr}, forms from decomposing
\SI{125}{\gram} of a 50.0\%-by-mass \ce{H2O2} solution?
\workspace[4.5cm]
\keyonly{\begin{apcanswer}mass \ce{H2O2} $=125\times0.500=62.5$~g,
$n=62.5/34.02=1.837$~mol

2:1 ratio with \ce{O2}:

\[n(\ce{O2}) = 1.837/2 = 0.9186\ \text{mol}\]

$P=746/760=0.9816$~atm, $T=300$~K:

\[V = \frac{nRT}{P} = \mathbf{23.0}\ \text{L}\]\end{apcanswer}}

\begin{rubric}
  \rpoint{1}{moles of \ce{H2O2} from the 50\% solution}
  \rpoint{1}{applies the 2:1 ratio to find $n(\ce{O2})$}
  \rpoint{1}{\SI{23.0}{\liter}}
\end{rubric}

% ---------------------------------------------------------------------------
\frq{long}{4}{Zumdahl \S5.4 \textbullet\ large-scale gas stoichiometry
  \emph{(after Zumdahl Q75)}}

In 1897, Swedish explorer Andr{\'e}e attempted a hydrogen-balloon flight
toward the North Pole; the \ce{H2} was generated from iron splints and
dilute sulfuric acid: \ce{Fe(s) + H2SO4(aq) -> FeSO4(aq) + H2(g)}. The
balloon's volume was $4800$~m$^3$, with an estimated 20.\% hydrogen loss
during filling.

Find the mass of iron splints and of 98\% (by mass) \ce{H2SO4} needed to
completely fill the balloon, assuming \SI{0}{\celsius}, \SI{1.0}{\atm}
during filling, 100\% yield.
\workspace[6.5cm]
\keyonly{\begin{apcanswer}\textbf{Step 1 --- moles of \ce{H2} needed in
the finished balloon:}

\[n = \frac{PV}{RT} = \frac{(1.0)(4.80\times10^6\ \text{L})}
{(0.08206)(273.15)} = 2.141\times10^5\ \text{mol}\]

\textbf{Step 2 --- account for the 20\% loss} --- 25\% more must be
\emph{produced} than ends up in the balloon:

\[n_{\text{produced}} = \frac{2.141\times10^5}{0.80}
= 2.677\times10^5\ \text{mol}\]

\textbf{Step 3 --- 1:1:1 stoichiometry} with \ce{Fe} and \ce{H2SO4}:

\[\text{mass Fe} = 2.677\times10^5\times55.85
= \mathbf{1.50\times10^4}\ \text{kg}\]

\[\text{mass pure \ce{H2SO4}} = 2.677\times10^5\times98.09
= 2.626\times10^7\ \text{g}\]

\[\text{mass 98\% solution} = \frac{2.626\times10^7}{0.98}
= \mathbf{2.68\times10^4}\ \text{kg}\]

(Metric-ton-scale masses are correct here, not a units slip --- this was a
genuinely enormous balloon.)\end{apcanswer}}

\begin{rubric}
  \rpoint{1}{moles of \ce{H2} needed in the final balloon}
  \rpoint{1}{correctly scales up for the 20\% loss}
  \rpoint{1}{mass of Fe, $\approx1.50\times10^4$~kg}
  \rpoint{1}{mass of 98\% \ce{H2SO4} solution, $\approx2.68\times10^4$~kg}
\end{rubric}

% ---------------------------------------------------------------------------
\frq{short}{3}{Zumdahl \S5.4 \textbullet\ two-step gas stoichiometry
  \emph{(after Zumdahl Q76)}}

\ce{S(s) + O2(g) -> SO2(g)}, then \ce{2SO2(g) + O2(g) -> 2SO3(g)}.

What volume of \ce{O2(g)} at \SI{350.}{\celsius} and \SI{5.25}{\atm} is
needed to fully convert \SI{5.00}{\gram} sulfur to \ce{SO3}?
\workspace[4.5cm]
\keyonly{\begin{apcanswer}$n(\ce{S})=5.00/32.07=0.1559$~mol. Each mole of
\ce{S} needs 1~mol \ce{O2} in step 1 and 0.5~mol \ce{O2} in step 2 ---
1.5~mol \ce{O2} total per mole \ce{S}:

\[n(\ce{O2}) = 0.1559\times1.5 = 0.2339\ \text{mol}\]

$T=623$~K:

\[V = \frac{nRT}{P} = \frac{(0.2339)(0.08206)(623)}{5.25}
= \mathbf{2.28}\ \text{L}\]\end{apcanswer}}

\begin{rubric}
  \rpoint{1}{moles of sulfur}
  \rpoint{1}{combines both steps into 1.5~mol \ce{O2} per mol \ce{S}}
  \rpoint{1}{\SI{2.28}{\liter}}
\end{rubric}

% ---------------------------------------------------------------------------
\frq{long}{4}{Zumdahl \S5.4 \textbullet\ limiting reactant, partial
  pressures given \emph{(after Zumdahl Q77)}}

A \SI{15.0}{\liter} rigid container holds \SI{0.500}{\atm} \ce{Kr} and
\SI{1.50}{\atm} \ce{Cl2} at \SI{350.}{\celsius}; they react:
\ce{Kr(g) + 2Cl2(g) -> KrCl4(s)}.

Find the mass of \ce{KrCl4} produced, assuming 100\% yield.
\workspace[5.5cm]
\keyonly{\begin{apcanswer}$T=623$~K

\[n(\ce{Kr}) = \frac{(0.500)(15.0)}{(0.08206)(623)} = 0.1467\ \text{mol}\]

\[n(\ce{Cl2}) = \frac{(1.50)(15.0)}{(0.08206)(623)} = 0.4401\ \text{mol}\]

Fully consuming the \ce{Kr} needs $2\times0.1467=0.2934$~mol \ce{Cl2}, and
\SI{0.4401}{\mole} is available --- \textbf{\ce{Kr} is limiting}.

\[\text{mass}(\ce{KrCl4}) = 0.1467\times225.60
= \mathbf{33.1}\ \text{g}\]\end{apcanswer}}

\begin{rubric}
  \rpoint{1}{$n(\ce{Kr})$ and $n(\ce{Cl2})$ from their partial pressures}
  \rpoint{1}{identifies \ce{Kr} as limiting, with the comparison shown}
  \rpoint{1}{correct molar mass of \ce{KrCl4}}
  \rpoint{1}{\SI{33.1}{\gram}}
\end{rubric}

% ---------------------------------------------------------------------------
\frq{long}{4}{Zumdahl \S5.4 \textbullet\ three-reactant limiting case
  \emph{(after Zumdahl Q78)}}

Acrylonitrile synthesis: \ce{2C3H6(g) + 2NH3(g) + 3O2(g) -> 2C3H3N(g)
+ 6H2O(g)}. A \SI{150.}{\liter} reactor at \SI{25}{\celsius} is charged to
partial pressures $P(\ce{C3H6})=\SI{0.500}{\mega\pascal}$,
$P(\ce{NH3})=\SI{0.800}{\mega\pascal}$,
$P(\ce{O2})=\SI{1.500}{\mega\pascal}$.

Find the mass of acrylonitrile (\ce{C3H3N}) obtainable.
\workspace[6.5cm]
\keyonly{\begin{apcanswer}Convert each pressure to atm ($\div101325$~Pa,
$\times10^6$) and find moles ($T=298$~K):

$n(\ce{C3H6})=30.27$~mol, $n(\ce{NH3})=48.43$~mol,
$n(\ce{O2})=90.81$~mol

Each reactant's maximum \ce{C3H3N} yield (using its own coefficient, all
1:1 with \ce{C3H6}/\ce{NH3}, 2:3 with \ce{O2}):
\ce{C3H6}$\to30.27$~mol, \ce{NH3}$\to48.43$~mol,
\ce{O2}$\to90.81\times\tfrac23=60.54$~mol.

\textbf{\ce{C3H6} is limiting}, despite \ce{O2} having by far the highest
partial pressure --- pressure alone does not determine the limiting
reactant, the coefficient-scaled yield does.

\[\text{mass}(\ce{C3H3N}) = 30.27\times53.06
= \mathbf{1610}\ \text{g}\]\end{apcanswer}}

\begin{rubric}
  \rpoint{1}{converts all three partial pressures to moles correctly}
  \rpoint{1}{computes each reactant's coefficient-scaled yield}
  \rpoint{1}{identifies \ce{C3H6} as limiting despite \ce{O2}'s higher
    pressure}
  \rpoint{1}{\SI{1610}{\gram} (1.61~kg)}
\end{rubric}

% ---------------------------------------------------------------------------
\frq{long}{4}{Zumdahl \S5.4 \textbullet\ liquid + gas limiting reactant
  \emph{(after Zumdahl Q79)}}

\SI{50.0}{\milli\liter} liquid methanol, \ce{CH3OH} (density
\SI{0.850}{\gram\per\milli\liter}), reacts with \SI{22.8}{\liter} \ce{O2}
at \SI{27}{\celsius}, \SI{2.00}{\atm}: \ce{2CH3OH(l) + 3O2(g) -> 2CO2(g)
+ 4H2O(g)}.

Find the moles of \ce{H2O} formed if the reaction completes.
\workspace[5.5cm]
\keyonly{\begin{apcanswer}mass \ce{CH3OH} $=50.0\times0.850=42.5$~g,
$n=42.5/32.04=1.326$~mol

\[n(\ce{O2}) = \frac{(2.00)(22.8)}{(0.08206)(300)} = 1.852\ \text{mol}\]

Fully consuming the methanol needs $1.5\times1.326=1.990$~mol \ce{O2},
but only \SI{1.852}{\mole} is available --- \textbf{\ce{O2} is limiting}.

\[n(\ce{H2O}) = 1.852\times\frac43 = \mathbf{2.47}\ \text{mol}\]
\end{apcanswer}}

\begin{rubric}
  \rpoint{1}{converts the liquid methanol volume to moles}
  \rpoint{1}{moles of \ce{O2} via the ideal gas law}
  \rpoint{1}{identifies \ce{O2} as limiting, with the comparison shown}
  \rpoint{1}{\SI{2.47}{\mole} \ce{H2O}}
\end{rubric}

% ---------------------------------------------------------------------------
\frq{long}{4}{Zumdahl \S5.4 \textbullet\ continuous-flow limiting reactant
  \emph{(after Zumdahl Q80)}}

Urea synthesis: \ce{2NH3(g) + CO2(g) -> H2NCONH2(s) + H2O(g)}. \ce{NH3} at
\SI{223}{\celsius}, \SI{90.}{\atm} flows into a reactor at
\SI{500.}{\liter\per\minute}; \ce{CO2} at \SI{223}{\celsius}, \SI{45}{\atm}
flows in at \SI{600.}{\liter\per\minute}.

Find the mass of urea produced per minute, assuming 100\% yield.
\workspace[6cm]
\keyonly{\begin{apcanswer}$T=496$~K

\[\frac{n(\ce{NH3})}{\text{min}} = \frac{(90.)(500.)}{(0.08206)(496)}
= 1106\ \text{mol/min}\]

\[\frac{n(\ce{CO2})}{\text{min}} = \frac{(45)(600.)}{(0.08206)(496)}
= 663.4\ \text{mol/min}\]

Fully consuming the \ce{NH3} needs $1106/2=553$~mol/min \ce{CO2}, and
\SI{663.4}{\mole\per\minute} is available --- \textbf{\ce{NH3} is
limiting}.

\[\frac{\text{mass urea}}{\text{min}} = \frac{1106}{2}\times60.06
= \mathbf{3.32\times10^4}\ \text{g/min}\ (33.2\ \text{kg/min})
\]\end{apcanswer}}

\begin{rubric}
  \rpoint{1}{moles per minute of \ce{NH3} from its flow rate}
  \rpoint{1}{moles per minute of \ce{CO2} from its flow rate}
  \rpoint{1}{identifies \ce{NH3} as limiting, with the comparison shown}
  \rpoint{1}{\SI{3.32e4}{\gram\per\minute} urea}
\end{rubric}

% ---------------------------------------------------------------------------
\frq{long}{4}{Zumdahl \S5.4 \textbullet\ balancing plus volume-based
  limiting reactant \emph{(after Zumdahl Q81)}}

\ce{HCN} is produced commercially from \ce{CH4(g)}, \ce{NH3(g)}, and
\ce{O2(g)} at high temperature; gaseous water is the other product.

\begin{parts}
  \item Write the balanced chemical equation.
        \workspace[2.5cm]
        \keyonly{\begin{apcanswer}\ce{2CH4(g) + 2NH3(g) + 3O2(g) -> 2HCN(g)
        + 6H2O(g)}\end{apcanswer}}
  \item Find the volume of \ce{HCN(g)} obtainable from \SI{20.0}{\liter}
        \ce{CH4(g)}, \SI{20.0}{\liter} \ce{NH3(g)}, and \SI{20.0}{\liter}
        \ce{O2(g)} --- all at the same temperature and pressure.
        \workspace[4.5cm]
        \keyonly{\begin{apcanswer}Same $T,P$ means volume tracks moles.
        Divide each given volume by its coefficient:
        \ce{CH4}: $20.0/2=10.00$; \ce{NH3}: $20.0/2=10.00$; \ce{O2}:
        $20.0/3=6.667$ --- \textbf{\ce{O2} is limiting}.

        \[V(\ce{HCN}) = 6.667\times2 = \mathbf{13.3}\ \text{L}\]
        \end{apcanswer}}
\end{parts}

\begin{rubric}
  \rpoint{1}{(a) correctly balanced equation}
  \rpoint{1}{(b) divides each volume by its coefficient}
  \rpoint{1}{(b) identifies \ce{O2} as limiting}
  \rpoint{1}{(b) \SI{13.3}{\liter} \ce{HCN}}
\end{rubric}

% ---------------------------------------------------------------------------
\frq{long}{4}{Zumdahl \S5.4 \textbullet\ percent yield from flow rates
  \emph{(after Zumdahl Q82)}}

\ce{C2H4(g) + H2(g) ->[\text{catalyst}] C2H6(g)}. \ce{C2H4} enters a
catalytic reactor at \SI{25.0}{\atm}, \SI{300.}{\celsius},
\SI{1000.}{\liter\per\minute}; \ce{H2} enters at \SI{25.0}{\atm},
\SI{300.}{\celsius}, \SI{1500.}{\liter\per\minute}. If \SI{15.0}{\kilo
\gram} \ce{C2H6} is collected per minute, find the percent yield.
\workspace[6cm]
\keyonly{\begin{apcanswer}$T=573$~K

\[\frac{n(\ce{C2H4})}{\text{min}} = \frac{(25.0)(1000.)}
{(0.08206)(573)} = 531.7\ \text{mol/min}\]

\[\frac{n(\ce{H2})}{\text{min}} = \frac{(25.0)(1500.)}
{(0.08206)(573)} = 797.5\ \text{mol/min}\]

1:1 ratio, so \textbf{\ce{C2H4} is limiting}.

\[\text{theoretical mass} = 531.7\times30.07
= 1.599\times10^4\ \text{g/min}\]

\[\%\ \text{yield} = \frac{1.50\times10^4}{1.599\times10^4}\times100
= \mathbf{93.8\%}\]\end{apcanswer}}

\begin{rubric}
  \rpoint{1}{moles per minute of each reactant}
  \rpoint{1}{identifies \ce{C2H4} as limiting}
  \rpoint{1}{correct theoretical yield}
  \rpoint{1}{93.8\% yield}
\end{rubric}

% ---------------------------------------------------------------------------
\frq{short}{3}{Zumdahl \S5.4 \textbullet\ identifying a gas from density
  \emph{(after Zumdahl Q83)}}

An unknown diatomic gas has a density of \SI{3.164}{\gram\per\liter} at
STP.

Identify the gas.
\workspace[3.5cm]
\keyonly{\begin{apcanswer}\[M = d\times22.42 = 3.164\times22.42
= 70.94\ \text{g/mol}\]

Atomic mass $=70.94/2=35.47$~g/mol --- matches chlorine.

\[\textbf{Gas: \ce{Cl2}}\]\end{apcanswer}}

\begin{rubric}
  \rpoint{1}{finds the molar mass via $M=d\times22.42$}
  \rpoint{1}{divides by 2 for the atomic mass}
  \rpoint{1}{identifies \ce{Cl2}}
\end{rubric}

% ---------------------------------------------------------------------------
\frq{long}{4}{Zumdahl \S5.4 \textbullet\ empirical to molecular formula
  \emph{(after Zumdahl Q84)}}

A compound with empirical formula \ce{CHCl}: a \SI{256}{\milli\liter}
flask at \SI{373}{\kelvin}, \SI{750.}{\torr} holds \SI{0.800}{\gram} of
the gaseous compound.

Find the molecular formula.
\workspace[5.5cm]
\keyonly{\begin{apcanswer}$P=750/760=0.9868$~atm

\[n = \frac{PV}{RT} = \frac{(0.9868)(0.256)}{(0.08206)(373)}
= 8.254\times10^{-3}\ \text{mol}\]

\[M = \frac{0.800}{8.254\times10^{-3}} = 96.9\ \text{g/mol}\]

Empirical formula mass of \ce{CHCl} $=12.01+1.008+35.45=48.47$~g/mol.

\[\frac{96.9}{48.47} \approx 2\]

\[\textbf{Molecular formula: \ce{C2H2Cl2}}\]\end{apcanswer}}

\begin{rubric}
  \rpoint{1}{moles of gas via the ideal gas law}
  \rpoint{1}{molar mass $\approx\SI{96.9}{\gram\per\mole}$}
  \rpoint{1}{empirical formula mass of \ce{CHCl}}
  \rpoint{1}{\ce{C2H2Cl2}, correctly justified by the multiplier of 2}
\end{rubric}

% ---------------------------------------------------------------------------
\frq{short}{2}{Zumdahl \S5.4 \textbullet\ density of a heavy vapor
  \emph{(after Zumdahl Q85)}}

Uranium hexafluoride is solid at room temperature but boils at
\SI{56}{\celsius}.

Find its density at \SI{60.}{\celsius} and \SI{745}{\torr}.
\workspace[3cm]
\keyonly{\begin{apcanswer}$M(\ce{UF6})=238.03+6(19.00)=352.03$~g/mol,
$P=0.9803$~atm, $T=333$~K

\[d = \frac{PM}{RT} = \mathbf{12.6}\ \text{g/L}\]\end{apcanswer}}

\begin{rubric}
  \rpoint{1}{molar mass of \ce{UF6} and correct unit conversions}
  \rpoint{1}{\SI{12.6}{\gram\per\liter}}
\end{rubric}

% ---------------------------------------------------------------------------
\frq{short}{3}{Zumdahl \S5.4 \textbullet\ comparing two vapor densities
  \emph{(after Zumdahl Q86)}}

Silicon tetrachloride (\ce{SiCl4}) and trichlorosilane (\ce{SiHCl3}) are
electronics-grade-silicon feedstocks.

Find the densities of pure \ce{SiCl4} vapor and pure \ce{SiHCl3} vapor at
\SI{85}{\celsius}, \SI{635}{\torr}.
\workspace[4.5cm]
\keyonly{\begin{apcanswer}$P=635/760=0.8355$~atm, $T=358$~K

$M(\ce{SiCl4})=28.09+4(35.45)=169.89$~g/mol:

\[d = \frac{PM}{RT} = \mathbf{4.83}\ \text{g/L}\]

$M(\ce{SiHCl3})=28.09+1.008+3(35.45)=135.45$~g/mol:

\[d = \frac{PM}{RT} = \mathbf{3.85}\ \text{g/L}\]\end{apcanswer}}

\begin{rubric}
  \rpoint{1}{molar masses of both compounds}
  \rpoint{1}{\ce{SiCl4}: \SI{4.83}{\gram\per\liter}}
  \rpoint{1}{\ce{SiHCl3}: \SI{3.85}{\gram\per\liter}}
\end{rubric}

% ============================================================================
% SESSION 4
% ============================================================================
\clearpage
\examsection{E}{Partial Pressure and Atmospheric Chemistry}
  {19 questions \textbullet\ 63 points}{\calculatorok}

% ---------------------------------------------------------------------------
\frq{short}{3}{Zumdahl \S5.5 \textbullet\ partial pressure from particle
  counts \emph{(after Zumdahl Q87)}}

A rigid container under a total pressure of \SI{1}{\atm} holds a mixture
of three noble gases. A particulate diagram shows 5 \ce{He} particles,
3 \ce{Ne} particles, and 2 \ce{Ar} particles (10 particles total).

Find the partial pressure of each gas.
\workspace[3.5cm]
\keyonly{\begin{apcanswer}Mole fraction equals particle-count fraction:

\[P(\ce{He}) = \tfrac{5}{10}(1) = \mathbf{0.500}\ \text{atm}\]
\[P(\ce{Ne}) = \tfrac{3}{10}(1) = \mathbf{0.300}\ \text{atm}\]
\[P(\ce{Ar}) = \tfrac{2}{10}(1) = \mathbf{0.200}\ \text{atm}\]\end{apcanswer}}

\begin{rubric}
  \rpoint{1}{recognizes particle-count fraction equals mole fraction}
  \rpoint{1}{all three partial pressures correct}
  \rpoint{1}{pressures sum to \SI{1.000}{\atm}}
\end{rubric}

% ---------------------------------------------------------------------------
\frq{long}{4}{Zumdahl \S5.5 \textbullet\ mixing two flasks
  \emph{(after Zumdahl Q88)}}

Two equal-volume flasks, same temperature, are joined by a closed
stopcock. The left flask holds only \ce{He}; a particulate diagram shows
5 \ce{He} particles. The right flask holds only \ce{Ne}; its diagram
shows 4 \ce{Ne} particles.

\begin{parts}
  \item Which is greater, the initial pressure of \ce{He} or \ce{Ne}? By
        how much?
        \workspace[3.5cm]
        \keyonly{\begin{apcanswer}$P\propto n$ at fixed $V,T$, so
        $P(\ce{He}):P(\ce{Ne})=5:4$ --- \ce{He}'s pressure is
        \textbf{25\% greater}.\end{apcanswer}}
  \item Once the stopcock opens, what volume does each gas occupy?
        \workspace[2.5cm]
        \keyonly{\begin{apcanswer}Both gases mix and spread through the
        \emph{combined} volume of both flasks.\end{apcanswer}}
  \item Express the final total pressure in terms of the original
        pressures $P(\ce{He})_i$ and $P(\ce{Ne})_i$.
        \workspace[3cm]
        \keyonly{\begin{apcanswer}Each gas's own volume doubles at
        constant $n,T$, so its own pressure halves:

        \[P_{\text{total,final}} = \frac{P(\ce{He})_i}{2}
        +\frac{P(\ce{Ne})_i}{2}
        = \mathbf{\dfrac{P(\ce{He})_i+P(\ce{Ne})_i}{2}}\]\end{apcanswer}}
  \item Express the final partial pressure of each gas in terms of its
        own original pressure.
        \workspace[3.5cm]
        \keyonly{\begin{apcanswer}$P(\ce{He})_f=\dfrac{P(\ce{He})_i}{2}$,\quad
        $P(\ce{Ne})_f=\dfrac{P(\ce{Ne})_i}{2}$ --- each gas simply
        halves, independently of the other.\end{apcanswer}}
\end{parts}

\begin{rubric}
  \rpoint{1}{(a) $P(\ce{He})$ is 25\% greater, from the 5:4 particle
    ratio}
  \rpoint{1}{(b) both gases occupy the full combined volume}
  \rpoint{1}{(c) $P_{\text{total,final}}=[P(\ce{He})_i+P(\ce{Ne})_i]/2$}
  \rpoint{1}{(d) each gas's final pressure is half its own original
    pressure}
\end{rubric}

% ---------------------------------------------------------------------------
\frq{short}{3}{Zumdahl \S5.5 \textbullet\ Dalton's law, masses given
  \emph{(after Zumdahl Q89)}}

For a deep dive, \SI{15.2}{\gram} \ce{He(g)} and \SI{30.6}{\gram}
\ce{O2(g)} are added to a previously evacuated \SI{5.00}{\liter} tank at
\SI{22}{\celsius}.

Find the partial pressure of each gas, and the total pressure.
\workspace[4.5cm]
\keyonly{\begin{apcanswer}$n(\ce{He})=15.2/4.003=3.797$~mol,
$n(\ce{O2})=30.6/32.00=0.9563$~mol, $T=295$~K

\[P(\ce{He}) = \frac{nRT}{V} = \mathbf{18.4}\ \text{atm}\]

\[P(\ce{O2}) = \mathbf{4.63}\ \text{atm}\]

\[P_{\text{total}} = 18.4+4.63 = \mathbf{23.0}\ \text{atm}\]
\end{apcanswer}}

\begin{rubric}
  \rpoint{1}{moles of both gases}
  \rpoint{1}{both partial pressures correct}
  \rpoint{1}{$P_{\text{total}}=\SI{23.0}{\atm}$}
\end{rubric}

% ---------------------------------------------------------------------------
\frq{short}{3}{Zumdahl \S5.5 \textbullet\ Dalton's law, light gases
  \emph{(after Zumdahl Q90)}}

\SI{1.00}{\gram} \ce{H2} and \SI{1.00}{\gram} \ce{He} are placed in a
\SI{1.00}{\liter} container at \SI{27}{\celsius}.

Find the partial pressure of each gas, and the total pressure.
\workspace[4cm]
\keyonly{\begin{apcanswer}$n(\ce{H2})=1.00/2.016=0.4960$~mol,
$n(\ce{He})=1.00/4.003=0.2498$~mol, $T=300$~K

\[P(\ce{H2}) = \mathbf{12.2}\ \text{atm}, \quad
P(\ce{He}) = \mathbf{6.15}\ \text{atm}\]

\[P_{\text{total}} = \mathbf{18.4}\ \text{atm}\]\end{apcanswer}}

\begin{rubric}
  \rpoint{1}{moles of both gases}
  \rpoint{1}{both partial pressures correct}
  \rpoint{1}{$P_{\text{total}}=\SI{18.4}{\atm}$}
\end{rubric}

% ---------------------------------------------------------------------------
\frq{short}{3}{Zumdahl \S5.5 \textbullet\ combining two flasks by Boyle's
  law \emph{(after Zumdahl Q91)}}

\SI{2.00}{\liter} \ce{H2} at \SI{475}{\torr} and \SI{1.00}{\liter} \ce{N2}
at \SI{0.200}{\atm} sit in separate flasks joined by a stopcock. Once
opened, the combined volume is \SI{3.00}{\liter} (constant temperature).

Find the final partial pressure of each gas, and the total pressure, in
torr.
\workspace[4.5cm]
\keyonly{\begin{apcanswer}Each gas expands into the full \SI{3.00}{\liter}
independently (Boyle's law, $P_f=P_iV_i/V_f$):

\[P(\ce{H2})_f = 475\times\frac{2.00}{3.00} = \mathbf{317}\ \text{torr}\]

\[P(\ce{N2})_f = 152\times\frac{1.00}{3.00} = \mathbf{50.7}\ \text{torr}\]

\[P_{\text{total}} = 317+50.7 = \mathbf{367}\ \text{torr}\]
\end{apcanswer}}

\begin{rubric}
  \rpoint{1}{applies Boyle's law to each gas independently}
  \rpoint{1}{both final partial pressures correct}
  \rpoint{1}{$P_{\text{total}}=\SI{367}{\torr}$}
\end{rubric}

% ---------------------------------------------------------------------------
\frq{short}{3}{Zumdahl \S5.5 \textbullet\ working backward from total
  pressure \emph{(after Zumdahl Q92)}}

Same two-flask apparatus as the previous problem: \SI{2.00}{\liter} \ce{H2}
at \SI{360.}{\torr} and \SI{1.00}{\liter} \ce{N2} at an unknown pressure.
After the stopcock opens (combined volume \SI{3.00}{\liter}), the total
pressure is \SI{320.}{\torr}.

Find the initial pressure of \ce{N2} in the \SI{1.00}{\liter} flask.
\workspace[4.5cm]
\keyonly{\begin{apcanswer}\ce{H2}'s contribution to the final pressure:

\[P(\ce{H2})_f = 360.\times\frac{2.00}{3.00} = 240.\ \text{torr}\]

\ce{N2}'s final contribution, by difference:

\[P(\ce{N2})_f = 320.-240. = 80.0\ \text{torr}\]

Boyle's law backward to the original \SI{1.00}{\liter} flask:

\[P(\ce{N2})_i = 80.0\times\frac{3.00}{1.00} = \mathbf{240.}\ \text{torr}\]
\end{apcanswer}}

\begin{rubric}
  \rpoint{1}{finds \ce{H2}'s final contribution, \SI{240.}{\torr}}
  \rpoint{1}{finds \ce{N2}'s final contribution by difference}
  \rpoint{1}{\SI{240.}{\torr} for \ce{N2}'s original pressure}
\end{rubric}

% ---------------------------------------------------------------------------
\frq{long}{4}{Zumdahl \S5.5 \textbullet\ three gases from mixed units
  \emph{(after Zumdahl Q94)}}

At \SI{0}{\celsius}, a \SI{1.0}{\liter} flask contains
$5.0\times10^{-2}$~mol \ce{N2}, $1.5\times10^2$~mg \ce{O2}, and
$5.0\times10^{21}$ molecules \ce{NH3}.

Find the partial pressure of each gas, and the total pressure.
\workspace[5.5cm]
\keyonly{\begin{apcanswer}$n(\ce{N2})=0.050$~mol,
$n(\ce{O2})=0.150/32.00=4.688\times10^{-3}$~mol,
$n(\ce{NH3})=5.0\times10^{21}/6.022\times10^{23}=8.303\times10^{-3}$~mol,
$T=273$~K

\[P(\ce{N2}) = \mathbf{1.12}\ \text{atm}, \quad
P(\ce{O2}) = \mathbf{0.105}\ \text{atm}, \quad
P(\ce{NH3}) = \mathbf{0.186}\ \text{atm}\]

\[P_{\text{total}} = \mathbf{1.41}\ \text{atm}\]\end{apcanswer}}

\begin{rubric}
  \rpoint{1}{converts mg \ce{O2} and molecules \ce{NH3} to moles correctly}
  \rpoint{1}{$P(\ce{N2})=\SI{1.12}{\atm}$}
  \rpoint{1}{$P(\ce{O2})$ and $P(\ce{NH3})$ both correct}
  \rpoint{1}{$P_{\text{total}}=\SI{1.41}{\atm}$}
\end{rubric}

% ---------------------------------------------------------------------------
\frq{short}{2}{Zumdahl \S5.5 \textbullet\ equimolar mixture
  \emph{(after Zumdahl Q96)}}

A 1:1 molar mixture of nitrous oxide and oxygen (a dental sedative) totals
\SI{2.50}{\atm}.

Find the partial pressure of each gas.
\workspace[2.5cm]
\keyonly{\begin{apcanswer}Equal moles $\Rightarrow$ equal partial
pressures:

\[P(\ce{N2O}) = P(\ce{O2}) = \mathbf{1.25}\ \text{atm each}\]
\end{apcanswer}}

\begin{rubric}
  \rpoint{1}{recognizes equal moles give equal partial pressures}
  \rpoint{1}{\SI{1.25}{\atm} each}
\end{rubric}

% ---------------------------------------------------------------------------
\frq{short}{3}{Zumdahl \S5.5 \textbullet\ partial pressure from mass
  fractions \emph{(after Zumdahl Q98)}}

A tank holds \SI{52.5}{\gram} \ce{O2} and \SI{65.1}{\gram} \ce{CO2} at
\SI{27}{\celsius}, total pressure \SI{9.21}{\atm}.

Find the partial pressure of each gas.
\workspace[4cm]
\keyonly{\begin{apcanswer}$n(\ce{O2})=52.5/32.00=1.641$~mol,
$n(\ce{CO2})=65.1/44.01=1.479$~mol

\[\chi(\ce{O2}) = \frac{1.641}{1.641+1.479} = 0.5259\]

\[P(\ce{O2}) = 0.5259\times9.21 = \mathbf{4.84}\ \text{atm}\]

\[P(\ce{CO2}) = 9.21-4.84 = \mathbf{4.37}\ \text{atm}\]\end{apcanswer}}

\begin{rubric}
  \rpoint{1}{moles of both gases from their masses}
  \rpoint{1}{mole fraction of \ce{O2}}
  \rpoint{1}{both partial pressures correct}
\end{rubric}

% ---------------------------------------------------------------------------
\frq{short}{3}{Zumdahl \S5.5 \textbullet\ collection over water, find
  volume \emph{(after Zumdahl Q100)}}

\ce{He} is collected over water at \SI{25}{\celsius}, \SI{1.00}{\atm}
total pressure. Water's vapor pressure at \SI{25}{\celsius} is
\SI{23.8}{\torr}.

What total gas volume must be collected to obtain \SI{0.586}{\gram}
\ce{He}?
\workspace[4.5cm]
\keyonly{\begin{apcanswer}$n(\ce{He})=0.586/4.003=0.1464$~mol

\[P(\ce{He}) = 760-23.8 = 736.2\ \text{torr} = 0.9687\ \text{atm}\]

\[V = \frac{nRT}{P(\ce{He})} = \frac{(0.1464)(0.08206)(298)}{0.9687}
= \mathbf{3.70}\ \text{L}\]

(This equals the \emph{total} collected volume, since \ce{He} and water
vapor share the same container.)\end{apcanswer}}

\begin{rubric}
  \rpoint{1}{moles of \ce{He}}
  \rpoint{1}{subtracts water's vapor pressure to find $P(\ce{He})$}
  \rpoint{1}{\SI{3.70}{\liter}}
\end{rubric}

% ---------------------------------------------------------------------------
\frq{long}{5}{Zumdahl \S5.5 \textbullet\ determining a product formula
  from pressure changes \emph{(after Zumdahl Q102)}}

A \SI{100.0}{\milli\liter} nickel container is filled with \ce{Xe} and
\ce{F2} at partial pressures \SI{1.24}{\atm} and \SI{10.10}{\atm}
respectively (\SI{25}{\celsius}). Heated to \SI{400}{\celsius}, the two
react to form a xenon fluoride; the vessel is then cooled so that leftover
\ce{F2} stays gaseous while the product becomes a nonvolatile solid. The
remaining \ce{F2} is transferred to a second \SI{100.0}{\milli\liter}
nickel container, where its pressure at \SI{25}{\celsius} is
\SI{7.62}{\atm}.

Assuming all the xenon reacted, find the formula of the product.
\workspace[6.5cm]
\keyonly{\begin{apcanswer}\textbf{Step 1 --- initial moles} (both at
298~K, \SI{0.1000}{\liter}):

\[n(\ce{Xe}) = \frac{(1.24)(0.1000)}{(0.08206)(298)}
= 5.071\times10^{-3}\ \text{mol}\]

\[n(\ce{F2})_i = \frac{(10.10)(0.1000)}{(0.08206)(298)}
= 4.130\times10^{-2}\ \text{mol}\]

\textbf{Step 2 --- remaining \ce{F2}} (second container, same $V,T$):

\[n(\ce{F2})_{\text{remaining}}
= \frac{(7.62)(0.1000)}{(0.08206)(298)} = 3.116\times10^{-2}\ \text{mol}\]

\textbf{Step 3 --- \ce{F2} consumed:}

\[n(\ce{F2})_{\text{consumed}} = 0.04130-0.03116
= 1.014\times10^{-2}\ \text{mol}\]

\textbf{Step 4 --- ratio to \ce{Xe}:}

\[\frac{n(\ce{F2})_{\text{consumed}}}{n(\ce{Xe})}
= \frac{1.014\times10^{-2}}{5.071\times10^{-3}} = \mathbf{2.00}\]

Each mole of \ce{F2} supplies 2~mol \ce{F} atoms, so a 2.00:1
\ce{F2}:\ce{Xe} ratio means 4~mol \ce{F} per mol \ce{Xe}:

\[\textbf{Product: \ce{XeF4}}\]\end{apcanswer}}

\begin{rubric}
  \rpoint{1}{initial moles of \ce{Xe} and \ce{F2}}
  \rpoint{1}{moles of \ce{F2} remaining, from the second container}
  \rpoint{1}{moles of \ce{F2} consumed, by subtraction}
  \rpoint{1}{correct \ce{F2}:\ce{Xe} ratio, $\approx2.00$}
  \rpoint{1}{converts the \ce{F2}:\ce{Xe} ratio to \ce{F}:\ce{Xe} and
    identifies \ce{XeF4}}
\end{rubric}

% ---------------------------------------------------------------------------
\frq{long}{4}{Zumdahl \S5.5 \textbullet\ percent yield, continuous flow
  \emph{(after Zumdahl Q103)}}

\ce{CO(g) + 2H2(g) -> CH3OH(g)}. \ce{H2} at STP flows into a reactor at
\SI{16.0}{\liter\per\minute}; \ce{CO} at STP flows in at
\SI{25.0}{\liter\per\minute}. If \SI{5.30}{\gram} methanol forms per
minute, find the percent yield.
\workspace[5.5cm]
\keyonly{\begin{apcanswer}\[\frac{n(\ce{H2})}{\text{min}}
= \frac{16.0}{22.42} = 0.7136\ \text{mol/min}, \quad
\frac{n(\ce{CO})}{\text{min}} = \frac{25.0}{22.42}
= 1.115\ \text{mol/min}\]

Need 2:1 \ce{H2}:\ce{CO}; available ratio is only 0.640:1 ---
\textbf{\ce{H2} is limiting}.

\[\text{theoretical mass} = \frac{0.7136}{2}\times32.04
= 11.43\ \text{g/min}\]

\[\%\ \text{yield} = \frac{5.30}{11.43}\times100
= \mathbf{46.4\%}\]\end{apcanswer}}

\begin{rubric}
  \rpoint{1}{moles per minute of both reactants}
  \rpoint{1}{identifies \ce{H2} as limiting}
  \rpoint{1}{correct theoretical yield}
  \rpoint{1}{46.4\% yield}
\end{rubric}

% ---------------------------------------------------------------------------
\frq{long}{4}{Zumdahl \S5.5 \textbullet\ fermentation headspace pressure
  \emph{(after Zumdahl Q104)}}

In the "M\'ethode Champenoise" process, grape juice ferments inside a
sealed bottle to produce sparkling wine:
\ce{C6H12O6(aq) -> 2C2H5OH(aq) + 2CO2(g)}. \SI{750.}{\milli\liter} grape
juice (density \SI{1.0}{\gram\per\milli\liter}) ferments in an
\SI{825}{\milli\liter} bottle until 12\% of the liquid's \emph{volume} is
ethanol (density \SI{0.79}{\gram\per\milli\liter}).

Assuming (incorrectly) that the \ce{CO2} is insoluble in the liquid, find
the \ce{CO2} pressure inside the bottle at \SI{25}{\celsius}.
\workspace[5.5cm]
\keyonly{\begin{apcanswer}ethanol volume $=0.12\times750=90.0$~mL,
mass $=90.0\times0.79=71.1$~g, $n=71.1/46.07=1.543$~mol

1:1 ratio with \ce{CO2}: $n(\ce{CO2})=1.543$~mol

Headspace volume $=825-750=75$~mL$=0.075$~L, $T=298$~K:

\[P(\ce{CO2}) = \frac{nRT}{V}
= \frac{(1.543)(0.08206)(298)}{0.075} = \mathbf{503}\ \text{atm}\]

This is not a typo --- it is a genuinely unphysical pressure (no ordinary
bottle survives \SI{500}{\atm}). That is the point of the exercise: the
assumption that \emph{all} the \ce{CO2} stays gaseous in a tiny headspace
is badly wrong, and the calculation shows exactly why real \ce{CO2}
\emph{must} stay dissolved until the cork is pulled.\end{apcanswer}}

\begin{rubric}
  \rpoint{1}{ethanol mass and moles from the 12\% volume condition}
  \rpoint{1}{applies the 1:1 mole ratio to find $n(\ce{CO2})$}
  \rpoint{1}{correct headspace volume, \SI{75}{\milli\liter}}
  \rpoint{1}{$\approx\SI{503}{\atm}$, with the physical interpretation}
\end{rubric}

% ---------------------------------------------------------------------------
\frq{short}{3}{Zumdahl \S5.5 \textbullet\ volume change on reaction
  \emph{(after Zumdahl Q106)}}

Equal moles of \ce{SO2} and \ce{O2} are mixed in a flexible (constant
$T,P$) container and sparked: \ce{2SO2(g) + O2(g) -> 2SO3(g)}.

Assuming the reaction completes, find the ratio of final to initial
gas-mixture volume.
\workspace[4cm]
\keyonly{\begin{apcanswer}Basis: $n$ mol \ce{SO2} + $n$ mol \ce{O2}
($2n$ total initially). The 2:1 ratio needs only $n/2$~mol \ce{O2} to
consume all the \ce{SO2} --- \textbf{\ce{SO2} is limiting}, \ce{O2} is in
excess.

\ce{O2} remaining $=n-n/2=n/2$; \ce{SO3} produced $=n$.

Final total $=n+n/2=1.5n$.

\[\frac{V_{\text{final}}}{V_{\text{initial}}} = \frac{1.5n}{2n}
= \mathbf{0.75}\ (3:4)\]\end{apcanswer}}

\begin{rubric}
  \rpoint{1}{identifies \ce{SO2} as limiting despite the equimolar setup}
  \rpoint{1}{correctly finds the final total moles, $1.5n$}
  \rpoint{1}{ratio $=0.75$ (or 3:4)}
\end{rubric}

% ---------------------------------------------------------------------------
\frq{long}{4}{Zumdahl \S5.5 \textbullet\ mixture composition from a
  pressure drop \emph{(after Zumdahl Q108)}}

Group 2A metal oxides react with \ce{CO2}: \ce{MO(s) + CO2(g) -> MCO3(s)}.
A \SI{2.85}{\gram} sample containing only \ce{MgO} and \ce{CuO} (only
\ce{MgO} reacts with \ce{CO2}) sits in a \SI{3.00}{\liter} container
filled with \ce{CO2} to \SI{740.}{\torr} at \SI{20.}{\celsius}. After the
reaction completes, the pressure is \SI{390.}{\torr} at the same
temperature.

Find the mass percent of \ce{MgO} in the mixture.
\workspace[5.5cm]
\keyonly{\begin{apcanswer}$T=293$~K

\[n(\ce{CO2})_i = \frac{(740/760)(3.00)}{(0.08206)(293)} = 0.1215\ \text{mol}\]

\[n(\ce{CO2})_f = \frac{(390/760)(3.00)}{(0.08206)(293)} = 0.06403\ \text{mol}\]

\[n(\ce{CO2})_{\text{consumed}} = n(\ce{MgO}) = 0.1215-0.06403
= 0.05747\ \text{mol}\]

$M(\ce{MgO})=24.31+16.00=40.31$~g/mol:

\[\text{mass \ce{MgO}} = 0.05747\times40.31 = 2.317\ \text{g}\]

\[\%\ \ce{MgO} = \frac{2.317}{2.85}\times100 = \mathbf{81.3\%}\]
\end{apcanswer}}

\begin{rubric}
  \rpoint{1}{moles of \ce{CO2} before and after the reaction}
  \rpoint{1}{moles of \ce{CO2} consumed equals moles of \ce{MgO}}
  \rpoint{1}{mass of \ce{MgO} using its correct molar mass}
  \rpoint{1}{81.3\%}
\end{rubric}

% ---------------------------------------------------------------------------
\frq{short}{2}{Zumdahl \S5.6 \textbullet\ atmospheric pressure and
  altitude \emph{(after Zumdahl Q126)}}

A \SI{1.0}{\liter} sample of air is collected at \SI{25}{\celsius}, sea
level (\SI{1.00}{\atm}).

Estimate this sample's volume at \SI{15}{\kilo\meter} altitude, where
pressure is about \SI{0.1}{\atm} (treat temperature as approximately
constant for this estimate).
\workspace[3cm]
\keyonly{\begin{apcanswer}Boyle's law:

\[V_2 = V_1\times\frac{P_1}{P_2} = 1.0\times\frac{1.00}{0.1}
= \mathbf{10}\ \text{L}\]\end{apcanswer}}

\begin{rubric}
  \rpoint{1}{applies Boyle's law}
  \rpoint{1}{\SI{10}{\liter}}
\end{rubric}

% ---------------------------------------------------------------------------
\frq{short}{2}{Zumdahl \S5.6 \textbullet\ acid rain
  \emph{(after Zumdahl Q128)}}

Write a balanced equation showing how sulfuric acid in acid rain reacts
with marble and limestone (both primarily calcium carbonate).
\workspace[2.5cm]
\keyonly{\begin{apcanswer}\[\ce{H2SO4(aq) + CaCO3(s) -> CaSO4(aq)
+ H2O(l) + CO2(g)}\]\end{apcanswer}}

\begin{rubric}
  \rpoint{1}{correct products (\ce{CaSO4}, \ce{H2O}, \ce{CO2})}
  \rpoint{1}{correctly balanced (all coefficients 1)}
\end{rubric}

% ---------------------------------------------------------------------------
\frq{long}{4}{Zumdahl \S5.6 \textbullet\ trace-gas concentration, ppmv
  \emph{(after Zumdahl Q129)}}

Atmospheric scientists report trace-compound concentrations in ppmv:
$\text{ppmv of X} = (V_X/V_{\text{air}})\times10^6$ at the same $T,P$
(equal to mole fraction $\times10^6$). On a November day, downtown
Denver's \ce{CO} concentration reached $3.0\times10^2$~ppmv; atmospheric
pressure was \SI{628}{\torr}, temperature \SI{0}{\celsius}.

\begin{parts}
  \item Find the partial pressure of \ce{CO}.
        \workspace[2.5cm]
        \keyonly{\begin{apcanswer}$\chi(\ce{CO})=3.0\times10^{-4}$;
        $P(\ce{CO})=3.0\times10^{-4}\times628=\mathbf{0.188}$~torr\end{apcanswer}}
  \item Find the \ce{CO} concentration in molecules per cubic meter.
        \workspace[3.5cm]
        \keyonly{\begin{apcanswer}$P(\ce{CO})=0.188/760
        =2.479\times10^{-4}$~atm, $T=273$~K

        \[\frac{n}{V} = \frac{P}{RT} = 1.107\times10^{-5}\ \text{mol/L}
        = 1.107\times10^{-2}\ \text{mol/m}^3\]

        \[\times 6.022\times10^{23}
        = \mathbf{6.66\times10^{21}}\ \text{molecules/m}^3\]
        \end{apcanswer}}
  \item Find the \ce{CO} concentration in molecules per cubic centimeter.
        \workspace[2.5cm]
        \keyonly{\begin{apcanswer}$6.66\times10^{21}\div10^6
        = \mathbf{6.66\times10^{15}}$~molecules/cm$^3$\end{apcanswer}}
\end{parts}

\begin{rubric}
  \rpoint{1}{(a) \SI{0.188}{\torr}}
  \rpoint{1}{(b) converts to molarity via $P/RT$}
  \rpoint{1}{(b) $6.66\times10^{21}$~molecules/m$^3$}
  \rpoint{1}{(c) $6.66\times10^{15}$~molecules/cm$^3$}
\end{rubric}

% ---------------------------------------------------------------------------
\frq{long}{4}{Zumdahl \S5.6 \textbullet\ two trace pollutants, ppmv
  \emph{(after Zumdahl Q130)}}

Trace organics are concentrated by pumping air through a porous trap,
then released for analysis by gas chromatography. A \SI{3.00}{\liter} air
sample at \SI{748}{\torr}, \SI{23}{\celsius} was passed through a trap,
yielding \SI{89.6}{\nano\gram} benzene (\ce{C6H6}) and
\SI{153}{\nano\gram} toluene (\ce{C7H8}).

Find the mixing ratio (ppmv, as in the previous problem) and the
concentration in molecules per cubic centimeter, for both compounds.
\workspace[6.5cm]
\keyonly{\begin{apcanswer}$P=748/760=0.9842$~atm, $T=296$~K

\[n_{\text{air}} = \frac{PV}{RT} = \frac{(0.9842)(3.00)}
{(0.08206)(296)} = 0.1216\ \text{mol}\]

\textbf{Benzene:} $n=89.6\times10^{-9}/78.11=1.147\times10^{-9}$~mol

\[\text{ppmv} = \frac{1.147\times10^{-9}}{0.1216}\times10^6
= \mathbf{9.44\times10^{-3}}\]

\[\text{conc.} = \frac{1.147\times10^{-9}\times6.022\times10^{23}}
{3000\ \text{cm}^3} = \mathbf{2.30\times10^{11}}\ \text{molecules/cm}^3\]

\textbf{Toluene:} $n=153\times10^{-9}/92.14=1.661\times10^{-9}$~mol

\[\text{ppmv} = \frac{1.661\times10^{-9}}{0.1216}\times10^6
= \mathbf{1.37\times10^{-2}}\]

\[\text{conc.} = \frac{1.661\times10^{-9}\times6.022\times10^{23}}
{3000\ \text{cm}^3} = \mathbf{3.33\times10^{11}}\ \text{molecules/cm}^3\]
\end{apcanswer}}

\begin{rubric}
  \rpoint{1}{total moles of air in the sampled volume}
  \rpoint{1}{both compounds' moles from their nanogram masses}
  \rpoint{1}{both ppmv values correct}
  \rpoint{1}{both molecules/cm$^3$ concentrations correct}
\end{rubric}

% ============================================================================
% SESSION 5
% ============================================================================
\clearpage
\examsection{F}{Kinetic Molecular Theory and Real Gases}
  {13 questions \textbullet\ 39 points}{\calculatorok}

% ---------------------------------------------------------------------------
\frq{short}{2}{Zumdahl \S5.6 \textbullet\ average kinetic energy
  \emph{(after Zumdahl Q109)}}

Find the average kinetic energy per mole of \ce{CH4(g)} \emph{and} of
\ce{N2(g)} at \SI{273}{\kelvin} and at \SI{546}{\kelvin}.
\workspace[3.5cm]
\keyonly{\begin{apcanswer}$KE=\tfrac32RT$ depends only on $T$, \textbf{not}
on molar mass --- so \ce{CH4} and \ce{N2} have the identical value at each
temperature:

\[KE(\SI{273}{\kelvin}) = \tfrac32(8.3145)(273)
= \mathbf{3.40\times10^3}\ \text{J/mol}\]

\[KE(\SI{546}{\kelvin}) = \mathbf{6.81\times10^3}\ \text{J/mol}\]
\end{apcanswer}}

\begin{rubric}
  \rpoint{1}{recognizes $KE$ is identical for both gases at a given $T$}
  \rpoint{1}{both correct values, and that doubling $T$ doubles $KE$}
\end{rubric}

% ---------------------------------------------------------------------------
\frq{short}{3}{Zumdahl \S5.6 \textbullet\ total kinetic energy of a
  mixture \emph{(after Zumdahl Q110)}}

A \SI{100.}{\liter} flask holds \ce{CH4} and \ce{Ar} at \SI{25}{\celsius}.
Argon present: \SI{228}{\gram}; methane's mole fraction is $0.650$.

Find the total kinetic energy of the mixture.
\workspace[4.5cm]
\keyonly{\begin{apcanswer}$n(\ce{Ar})=228/39.95=5.707$~mol.
$\chi(\ce{Ar})=0.350$, so:

\[n_{\text{total}} = \frac{5.707}{0.350} = 16.31\ \text{mol}\]

$KE_{\text{total}}$ needs only $n_{\text{total}}$, since $KE$ per mole is
identical for every gas at a given $T$:

\[KE_{\text{total}} = n_{\text{total}}\times\tfrac32RT
= 16.31\times1.5\times8.3145\times298
= \mathbf{6.06\times10^4}\ \text{J}\ (60.6\ \text{kJ})\]\end{apcanswer}}

\begin{rubric}
  \rpoint{1}{moles of \ce{Ar}}
  \rpoint{1}{total moles from the mole fraction}
  \rpoint{1}{\SI{60.6}{\kilo\joule}}
\end{rubric}

% ---------------------------------------------------------------------------
\frq{long}{4}{Zumdahl \S5.6 \textbullet\ four root-mean-square velocities
  \textbf{[beyond CED]} \emph{(after Zumdahl Q111)}}

Find the root-mean-square velocity of \ce{CH4(g)} and of \ce{N2(g)}, each
at \SI{273}{\kelvin} and at \SI{546}{\kelvin} (four values total).
\workspace[5cm]
\keyonly{\begin{apcanswer}$M(\ce{CH4})=\SI{0.01604}{\kilo\gram\per\mole}$,
$M(\ce{N2})=\SI{0.02802}{\kilo\gram\per\mole}$

\[\ce{CH4},\ \SI{273}{\kelvin}:\ \mathbf{652}\ \text{m/s} \qquad
\ce{CH4},\ \SI{546}{\kelvin}:\ \mathbf{921}\ \text{m/s}\]

\[\ce{N2},\ \SI{273}{\kelvin}:\ \mathbf{493}\ \text{m/s} \qquad
\ce{N2},\ \SI{546}{\kelvin}:\ \mathbf{697}\ \text{m/s}\]

At fixed $T$, the lighter gas (\ce{CH4}) always moves faster; doubling $T$
scales every speed by $\sqrt2$.\end{apcanswer}}

\begin{rubric}
  \rpoint{1}{correct formula $u_{\text{rms}}=\sqrt{3RT/M}$ applied with
    $M$ in kg/mol}
  \rpoint{1}{both \ce{CH4} values correct}
  \rpoint{1}{both \ce{N2} values correct}
  \rpoint{1}{notes the $\sqrt2$ scaling and that \ce{CH4} is always faster}
\end{rubric}

% ---------------------------------------------------------------------------
\frq{short}{2}{Zumdahl \S5.6 \textbullet\ matching speeds across two
  gases \textbf{[beyond CED]} \emph{(after Zumdahl Q112)}}

Separate \SI{1.0}{\liter} samples of \ce{He(g)} and \ce{UF6(g)}, both at
\SI{1.00}{\atm} with equal moles.

What temperature ratio $T(\ce{UF6})/T(\ce{He})$ gives them equal
root-mean-square velocity?
\workspace[3cm]
\keyonly{\begin{apcanswer}Setting $u_{\text{rms}}$ equal cancels the
common factors:

\[\frac{T(\ce{UF6})}{T(\ce{He})} = \frac{M(\ce{UF6})}{M(\ce{He})}
= \frac{352.03}{4.003} = \mathbf{87.9}\]

\ce{UF6} would need to be almost 88 times hotter (in kelvin) than
\ce{He} to move at the same speed.\end{apcanswer}}

\begin{rubric}
  \rpoint{1}{sets $\sqrt{3RT_{\ce{He}}/M_{\ce{He}}}
    =\sqrt{3RT_{\ce{UF6}}/M_{\ce{UF6}}}$ and simplifies}
  \rpoint{1}{ratio $\approx88$}
\end{rubric}

% ---------------------------------------------------------------------------
\frq{short}{2}{Zumdahl \S5.6 \textbullet\ conceptual, shrinking a piston
  \emph{(after Zumdahl Q113)}}

A gas-filled cylinder with a movable piston, held at \SI{1.00}{\atm},
shows its volume decreasing.

State two \emph{distinct} changes to the system that would produce this
volume decrease, and briefly explain why each works.
\workspace[3.5cm]
\keyonly{\begin{apcanswer}\textbf{(1)} Lower the temperature at constant
$n,P$ --- $V\propto T$, so cooling shrinks the volume.

\textbf{(2)} Remove some gas (decrease $n$) at constant $T,P$ ---
$V\propto n$, so less gas occupies less space.\end{apcanswer}}

\begin{rubric}
  \rpoint{1}{lowering temperature, with the $V\propto T$ reasoning}
  \rpoint{1}{removing gas, with the $V\propto n$ reasoning}
\end{rubric}

% ---------------------------------------------------------------------------
\frq{short}{3}{Zumdahl \S5.6 \textbullet\ conceptual, doubling a piston's
  volume \emph{(after Zumdahl Q114)}}

A cylinder with a weighted piston shows its volume doubling once the
weights are removed.

State three \emph{distinct} changes that would double the gas's volume,
and briefly explain each.
\workspace[4.5cm]
\keyonly{\begin{apcanswer}\textbf{(1)} Halve the pressure at constant
$n,T$ --- $V\propto1/P$ (this is exactly what removing weight from the
piston does).

\textbf{(2)} Double the absolute temperature at constant $n,P$ ---
$V\propto T$.

\textbf{(3)} Double the moles of gas at constant $T,P$ --- $V\propto n$.
\end{apcanswer}}

\begin{rubric}
  \rpoint{1}{halving pressure, with reasoning}
  \rpoint{1}{doubling temperature, with reasoning}
  \rpoint{1}{doubling moles, with reasoning}
\end{rubric}

% ---------------------------------------------------------------------------
\frq{long}{4}{Zumdahl \S5.6 \textbullet\ KMT reasoning across four changes
  \emph{(after Zumdahl Q115)}}

For a \SI{1.0}{\liter} container of neon gas at STP, state whether average
kinetic energy, average velocity, and wall-collision frequency increase,
decrease, or stay the same when:

\begin{parts}
  \item Temperature rises to \SI{100}{\celsius}.
        \workspace[2.5cm]
        \keyonly{\begin{apcanswer}All three \textbf{increase} --- faster molecules
        that also hit the walls more often.\end{apcanswer}}
  \item Temperature drops to \SI{-50}{\celsius}.
        \workspace[2.5cm]
        \keyonly{\begin{apcanswer}All three \textbf{decrease}.\end{apcanswer}}
  \item Volume drops to \SI{0.5}{\liter} ($T,n$ constant).
        \workspace[3.5cm]
        \keyonly{\begin{apcanswer}$KE$ and average velocity are \textbf{unchanged}
        (depend only on $T$); collision frequency \textbf{increases} ---
        the same particles now pack into half the space, so they strike
        the walls more often even at the same speed.\end{apcanswer}}
  \item The moles of neon double ($T,V$ constant).
        \workspace[3.5cm]
        \keyonly{\begin{apcanswer}$KE$ and average velocity are \textbf{unchanged};
        collision frequency \textbf{increases} --- twice as many
        particles means twice as many wall strikes per second.\end{apcanswer}}
\end{parts}

\begin{rubric}
  \rpoint{1}{(a) all three increase}
  \rpoint{1}{(b) all three decrease}
  \rpoint{1}{(c) $KE$/velocity unchanged, collision frequency increases,
    with reasoning}
  \rpoint{1}{(d) $KE$/velocity unchanged, collision frequency increases,
    with reasoning}
\end{rubric}

% ---------------------------------------------------------------------------
\frq{long}{4}{Zumdahl \S5.6 \textbullet\ comparing two unknown gases
  \emph{(after Zumdahl Q116)}}

Two \SI{1.0}{\liter} containers hold different gases at the same
temperature and pressure. Container A holds \SI{0.34}{\gram}; container B
holds \SI{0.48}{\gram}.

\begin{parts}
  \item Which sample has more molecules present? Explain.
        \workspace[2.5cm]
        \keyonly{\begin{apcanswer}\textbf{Equal} --- same $T,P,V$ means equal moles
        by $PV=nRT$, regardless of mass or identity.\end{apcanswer}}
  \item Which sample has the greater average kinetic energy? Explain.
        \workspace[2.5cm]
        \keyonly{\begin{apcanswer}\textbf{Equal} --- $KE$ depends only on $T$.\end{apcanswer}}
  \item Which sample has the faster average velocity? Explain.
        \workspace[3.5cm]
        \keyonly{\begin{apcanswer}\textbf{A} (the lighter gas) --- since
        $u_{\text{rms}}\propto1/\sqrt{M}$ at fixed $T$, and A has fewer
        grams for the same moles, A's molar mass is smaller.\end{apcanswer}}
  \item If B's heavier molecules hit the walls harder, how can the
        pressure in both containers be equal?
        \workspace[3.5cm]
        \keyonly{\begin{apcanswer}B's molecules move proportionally
        \emph{slower} (lower $u_{\text{rms}}$) by exactly the amount that
        offsets their greater mass per collision --- the two effects
        cancel, which is exactly what makes $PV=nRT$ identity-independent.\end{apcanswer}}
\end{parts}

\begin{rubric}
  \rpoint{1}{(a) equal molecules, justified via $PV=nRT$}
  \rpoint{1}{(b) equal $KE$, justified via $T$}
  \rpoint{1}{(c) A faster, justified via $u_{\text{rms}}\propto1/\sqrt M$}
  \rpoint{1}{(d) explains the mass/speed trade-off that keeps pressure
    equal}
\end{rubric}

% ---------------------------------------------------------------------------
\frq{long}{3}{Zumdahl \S5.6 \textbullet\ ranking four gases
  \textbf{[beyond CED]} \emph{(after Zumdahl Q118)}}

Separate \SI{1.0}{\liter} samples of \ce{H2}, \ce{Xe}, \ce{Cl2}, and
\ce{O2}, all at STP.

\begin{parts}
  \item Rank the gases by increasing average kinetic energy.
        \workspace[2.5cm]
        \keyonly{\begin{apcanswer}\textbf{All equal} --- $KE$ depends only on $T$,
        which is the same (STP) for all four.\end{apcanswer}}
  \item Rank the gases by increasing average velocity.
        \workspace[3.5cm]
        \keyonly{\begin{apcanswer}\ce{Xe} (228 m/s) $<$ \ce{Cl2} (310 m/s)
        $<$ \ce{O2} (461 m/s) $<$ \textbf{\ce{H2}} (1840 m/s) ---
        lightest moves fastest.\end{apcanswer}}
  \item How could \ce{O2} and \ce{H2} samples ever have the \emph{same}
        average velocity?
        \workspace[2.5cm]
        \keyonly{\begin{apcanswer}Only at different temperatures: \ce{O2} would
        need to be $M(\ce{O2})/M(\ce{H2})\approx16$ times hotter (in
        kelvin) than \ce{H2}.\end{apcanswer}}
\end{parts}

\begin{rubric}
  \rpoint{1}{(a) all equal, with the $T$-only reasoning}
  \rpoint{1}{(b) correct increasing order by $u_{\text{rms}}$}
  \rpoint{1}{(c) identifies that only a $\approx16\times$ temperature gap
    could equalize them}
\end{rubric}

% ---------------------------------------------------------------------------
\frq{short}{3}{Zumdahl \S5.7 \textbullet\ identifying a gas by effusion
  \textbf{[beyond CED]} \emph{(after Zumdahl Q119)}}

An unknown gas effuses at \SI{31.50}{\milli\liter\per\minute}; \ce{O2(g)}
effuses at \SI{30.50}{\milli\liter\per\minute} under identical conditions.

Identify the unknown gas (candidates: \ce{CH4}, \ce{CO}, \ce{NO},
\ce{CO2}, \ce{NO2}).
\workspace[4cm]
\keyonly{\begin{apcanswer}\[M_{\text{unknown}} = M(\ce{O2})\times
\left(\frac{\text{rate}(\ce{O2})}{\text{rate}_{\text{unknown}}}\right)^2
= 32.00\times\left(\frac{30.50}{31.50}\right)^2
= \mathbf{30.0}\ \text{g/mol}\]

This matches \ce{NO} ($\SI{30.01}{\gram\per\mole}$) almost exactly ---
every other candidate is off by more than 1~g/mol.

\[\textbf{Identity: \ce{NO}}\]\end{apcanswer}}

\begin{rubric}
  \rpoint{1}{correct Graham's law setup, with \ce{O2}'s rate in the
    numerator}
  \rpoint{1}{$M\approx\SI{30.0}{\gram\per\mole}$}
  \rpoint{1}{identifies \ce{NO}, ruling out the other candidates}
\end{rubric}

% ---------------------------------------------------------------------------
\frq{short}{2}{Zumdahl \S5.7 \textbullet\ unknown molar mass from
  effusion \textbf{[beyond CED]} \emph{(after Zumdahl Q120)}}

An unknown gas effuses at \SI{24.0}{\milli\liter\per\minute}; pure
\ce{CH4(g)} effuses at \SI{47.8}{\milli\liter\per\minute} under identical
conditions.

Find the unknown gas's molar mass.
\workspace[3cm]
\keyonly{\begin{apcanswer}\[M_{\text{unknown}} = M(\ce{CH4})\times
\left(\frac{\text{rate}(\ce{CH4})}{\text{rate}_{\text{unknown}}}\right)^2
= 16.04\times\left(\frac{47.8}{24.0}\right)^2
= \mathbf{63.6}\ \text{g/mol}\]\end{apcanswer}}

\begin{rubric}
  \rpoint{1}{correct Graham's law setup}
  \rpoint{1}{\SI{63.6}{\gram\per\mole}}
\end{rubric}

% ---------------------------------------------------------------------------
\frq{short}{2}{Zumdahl \S5.7 \textbullet\ effusion time comparison
  \textbf{[beyond CED]} \emph{(after Zumdahl Q122)}}

\SI{1.0}{\liter} helium takes \SI{4.5}{\minute} to effuse through a porous
barrier.

How long would \SI{1.0}{\liter} \ce{Cl2} gas take, under identical
conditions?
\workspace[3cm]
\keyonly{\begin{apcanswer}Effusion \emph{time} scales as $\sqrt{M}$
(inverse of rate):

\[t(\ce{Cl2}) = t(\ce{He})\times\sqrt{\frac{M(\ce{Cl2})}{M(\ce{He})}}
= 4.5\times\sqrt{\frac{70.90}{4.003}} = \mathbf{18.9}\ \text{min}\]
\end{apcanswer}}

\begin{rubric}
  \rpoint{1}{recognizes time scales as $\sqrt M$, not $1/\sqrt M$}
  \rpoint{1}{\SI{18.9}{\minute}}
\end{rubric}

% ---------------------------------------------------------------------------
\frq{long}{5}{Zumdahl \S5.8 \textbullet\ van der Waals versus ideal
  \textbf{[beyond CED]} \emph{(after Zumdahl Q124)}}

\SI{0.5000}{\mole} \ce{N2} sits in a \SI{10.000}{\liter} container at
\SI{25.0}{\celsius}.

Van der Waals constants for \ce{N2}:
$a=\SI{1.39}{\liter\squared\atm\per\mole\squared}$,
$b=\SI{0.0391}{\liter\per\mole}$.

\begin{parts}
  \item Find the pressure using the ideal gas law.
        \workspace[2.5cm]
        \keyonly{\begin{apcanswer}$P=nRT/V=\mathbf{1.223}$~atm\end{apcanswer}}
  \item Find the pressure using the van der Waals equation.
        \workspace[3.5cm]
        \keyonly{\begin{apcanswer}\[P = \frac{nRT}{V-nb}
        - a\left(\frac{n}{V}\right)^2
        = \frac{(0.5000)(0.08206)(298.15)}{10.000-(0.5000)(0.0391)}
        - 1.39\left(\frac{0.5000}{10.000}\right)^2
        = \mathbf{1.222}\ \text{atm}\]\end{apcanswer}}
  \item Compare the two results.
        \workspace[2.5cm]
        \keyonly{\begin{apcanswer}They differ by only about
        \textbf{0.09\%} --- at \SI{10.0}{\liter} this gas behaves almost
        ideally.\end{apcanswer}}
  \item Repeat parts (a)--(b) for the \emph{same} \SI{0.5000}{\mole} in a
        \SI{1.0000}{\liter} container instead, and compare the size of the
        deviation to part (c).
        \workspace[4cm]
        \keyonly{\begin{apcanswer}Ideal: $P=nRT/V=\mathbf{12.23}$~atm.

        Van der Waals: $P=\dfrac{nRT}{V-nb}-a(n/V)^2=\mathbf{12.13}$~atm.

        \% difference $\approx\mathbf{0.85\%}$ --- about \textbf{10 times
        larger} than the 10.0~L case. Compressing the same moles into a
        tenth of the volume raises the pressure roughly tenfold, and the
        van der Waals corrections grow with it: this is exactly the
        "close to condensation / high pressure" regime CED~3.6 describes
        qualitatively.\end{apcanswer}}
\end{parts}

\begin{rubric}
  \rpoint{1}{(a) \SI{1.223}{\atm}}
  \rpoint{1}{(b) \SI{1.222}{\atm}}
  \rpoint{1}{(c) correctly compares, $\approx0.09\%$ difference}
  \rpoint{2}{(d) both 1.0~L values correct AND correctly identifies the
    much larger ($\approx10\times$) deviation at higher pressure}
\end{rubric}

% ============================================================================
% SESSION 6
% ============================================================================
\clearpage
\examsection{G}{Challenge Extension}{13 questions \textbullet\ 75 points
  \textbullet\ separate sitting}{\calculatorok}

\begin{apcnote}[About Part G]
Every Challenge Problem printed on Zumdahl's Chapter~5 review pages is
here. These go well past what the AP exam will ask -- multi-gas mixture
mass balances, hot-air-balloon lift calculations, and van der Waals
derivations -- but every technique is one you already used in Sessions
1--5, just layered together.
\end{apcnote}

% ---------------------------------------------------------------------------
\frq{long}{6}{Zumdahl \S5.4 \textbullet\ two-oxide mixture from a pressure
  drop \emph{(after Zumdahl Q157)}}

A \SI{5.14}{\gram} mixture of \ce{BaO(s)} and \ce{CaO(s)} is placed in a
\SI{1.50}{\liter} flask containing \ce{CO2(g)} at \SI{30.0}{\celsius} and
\SI{750.}{\torr}: \ce{BaO(s) + CO2(g) -> BaCO3(s)} and
\ce{CaO(s) + CO2(g) -> CaCO3(s)}. After the reaction completes, the
remaining \ce{CO2(g)} pressure is \SI{230.}{\torr} at the same
temperature.

Find the mass percent of \ce{CaO} and of \ce{BaO} in the original mixture.
\workspace[7.5cm]
\keyonly{\begin{apcanswer}$T=303.15$~K

\[n(\ce{CO2})_i = \frac{(750/760)(1.50)}{(0.08206)(303.15)}
= 0.05950\ \text{mol}\]

\[n(\ce{CO2})_f = \frac{(230/760)(1.50)}{(0.08206)(303.15)}
= 0.01825\ \text{mol}\]

\[n(\ce{CO2})_{\text{consumed}} = 0.05950-0.01825 = 0.04126\ \text{mol}
= x(\ce{BaO})+x(\ce{CaO})\]

Both oxides react 1:1 with \ce{CO2}, so this is a system in two unknowns
($x=$~mol \ce{BaO}, $y=$~mol \ce{CaO}$)$:

\[x+y = 0.04126 \qquad 153.33x+56.08y = 5.14\]

Solving: $x=0.02906$~mol \ce{BaO}, $y=0.01219$~mol \ce{CaO}.

mass \ce{BaO} $=0.02906\times153.33=4.456$~g;
mass \ce{CaO} $=0.01219\times56.08=0.684$~g.

\[\%\ce{BaO} = \frac{4.456}{5.14}\times100 = \mathbf{86.7\%} \qquad
\%\ce{CaO} = \mathbf{13.3\%}\]\end{apcanswer}}

\begin{rubric}
  \rpoint{1}{moles of \ce{CO2} before the reaction}
  \rpoint{1}{moles of \ce{CO2} after the reaction}
  \rpoint{1}{moles of \ce{CO2} consumed equals total moles of oxide}
  \rpoint{1}{sets up the correct two-equation system (mass and mole
    balance)}
  \rpoint{1}{solves for both masses correctly}
  \rpoint{1}{86.7\% \ce{BaO}, 13.3\% \ce{CaO}}
\end{rubric}

% ---------------------------------------------------------------------------
\frq{long}{6}{Zumdahl \S5.4 \textbullet\ metal-alloy composition from
  \ce{H2} volume \emph{(after Zumdahl Q158)}}

A \SI{0.362}{\gram} mixture of chromium and zinc is reacted with excess
hydrochloric acid: \ce{Zn(s) + 2HCl(aq) -> ZnCl2(aq) + H2(g)} and
\ce{2Cr(s) + 6HCl(aq) -> 2CrCl3(aq) + 3H2(g)}. After all the metal
reacts, \SI{225}{\milli\liter} dry \ce{H2} is collected at
\SI{27}{\celsius}, \SI{750.}{\torr}.

Determine the mass percent of \ce{Zn} in the metal sample.
\workspace[7.5cm]
\keyonly{\begin{apcanswer}\[n(\ce{H2}) = \frac{(750/760)(0.225)}
{(0.08206)(300)} = 0.009019\ \text{mol}\]

Each mole \ce{Zn} produces 1~mol \ce{H2}; each mole \ce{Cr} produces
1.5~mol \ce{H2}. With $x=$~mol \ce{Zn}, $y=$~mol \ce{Cr}:

\[65.38x+52.00y = 0.362 \qquad x+1.5y = 0.009019\]

Solving: $y=0.004942$~mol \ce{Cr}, $x=0.001606$~mol \ce{Zn}.

mass \ce{Zn} $=0.001606\times65.38=0.1050$~g;
mass \ce{Cr} $=0.004942\times52.00=0.2570$~g (sum checks: 0.3620~g).

\[\%\ce{Zn} = \frac{0.1050}{0.362}\times100 = \mathbf{29.0\%}\]
\end{apcanswer}}

\begin{rubric}
  \rpoint{1}{moles of \ce{H2} collected}
  \rpoint{1}{correctly relates 1~mol \ce{Zn} to 1~mol \ce{H2} and 1~mol
    \ce{Cr} to 1.5~mol \ce{H2}}
  \rpoint{1}{sets up the two-equation system}
  \rpoint{1}{solves for both metals' moles}
  \rpoint{1}{both masses check against the \SI{0.362}{\gram} total}
  \rpoint{1}{29.0\% \ce{Zn}}
\end{rubric}

% ---------------------------------------------------------------------------
\frq{long}{6}{Zumdahl \S5.4 \textbullet\ identifying a hydrocarbon from
  its combustion products \emph{(after Zumdahl Q159)}}

A hydrocarbon sample (\ce{C} and \ce{H} only) at \SI{0.959}{\atm},
\SI{298}{\kelvin} is burned completely in oxygen, producing a mixture of
gaseous \ce{CO2} and \ce{H2O} at \SI{1.51}{\atm}, \SI{375}{\kelvin},
density \SI{1.391}{\gram\per\liter}, occupying a volume \textbf{four
times} as large as the original hydrocarbon sample's.

Determine the molecular formula of the hydrocarbon.
\workspace[8cm]
\keyonly{\begin{apcanswer}Let the hydrocarbon be $\ce{C}_x\ce{H}_y$, so

\[\ce{C}_x\ce{H}_y + \left(x+\frac{y}{4}\right)\ce{O2}
\longrightarrow x\,\ce{CO2} + \frac{y}{2}\ce{H2O}\]

\textbf{Step 1 --- total product moles per mole hydrocarbon.} With
$n_{\text{HC}}=P_1V/(RT_1)$ and $n_{\text{products}}=P_2(4V)/(RT_2)$, the
sample volume $V$ cancels:

\[\frac{n_{\text{products}}}{n_{\text{HC}}}
= \frac{4P_2T_1}{P_1T_2}
= \frac{4(1.51)(298)}{(0.959)(375)} = 5.00 = x+\frac{y}{2}\]

(this is exactly the moles of \ce{CO2} plus \ce{H2O} produced per mole of
hydrocarbon burned).

\textbf{Step 2 --- average molar mass of the product mixture,} from its
density:

\[M_{\text{avg}} = \frac{dRT_2}{P_2}
= \frac{(1.391)(0.08206)(375)}{1.51} = 28.3\ \text{g/mol}\]

\textbf{Step 3 --- solve the two equations.} With
$A=x+y/2=5.00$ and $M_{\text{avg}}=[44.01x+18.02(A-x)]/A$:

\[x = \frac{A(M_{\text{avg}}-18.02)}{44.01-18.02} \approx 2, \qquad
\frac{y}{2} = A-x \approx 3 \Rightarrow y\approx6\]

\[\textbf{Molecular formula: \ce{C2H6} (ethane)}\]

Check: pure \ce{C2H6} combustion gives $A=x+y/2=2+3=5$ exactly, and
$M_{\text{avg}}=[2(44.01)+3(18.02)]/5=28.4$~g/mol --- both match the
measured values to within rounding.\end{apcanswer}}

\begin{rubric}
  \rpoint{1}{writes the general combustion equation for $\ce{C}_x\ce{H}_y$}
  \rpoint{1}{finds $x+y/2\approx5.00$ from the mole-ratio relationship
    (volume must cancel)}
  \rpoint{1}{finds $M_{\text{avg}}\approx\SI{28.3}{\gram\per\mole}$ of the
    product mixture}
  \rpoint{2}{solves the two equations for $x\approx2$, $y\approx6$}
  \rpoint{1}{identifies \ce{C2H6} and verifies against both checks}
\end{rubric}

% ---------------------------------------------------------------------------
\frq{long}{7}{Zumdahl \S5.4 \textbullet\ mixture density before and after
  reaction \emph{(after Zumdahl Q160)}}

An equimolar mixture of \ce{SO2} and \ce{O2}, plus some \ce{He}, sits in a
piston-fitted container (constant $T,P$). The mixture's density at STP is
\SI{1.924}{\gram\per\liter}.

\begin{parts}
  \item What is the mole fraction of \ce{He} in the original mixture?
        \workspace[5cm]
        \keyonly{\begin{apcanswer}Let $m=$~mol \ce{SO2}$=$mol \ce{O2}
        (equimolar), $h=$~mol \ce{He}.

        \[M_{\text{avg}} = d\times22.42 = 1.924\times22.42
        = 43.14\ \text{g/mol}\]

        \[M_{\text{avg}}(2m+h) = m(64.07)+m(32.00)+h(4.003)\]

        Solving for the ratio $m/h$:

        \[\frac{m}{h} = \frac{4.003-M_{\text{avg}}}{2M_{\text{avg}}-96.07}
        = \mathbf{3.994}\]

        \[\chi(\ce{He}) = \frac{h}{2m+h} = \frac{1}{2(3.994)+1}
        = \mathbf{0.111}\ (11.1\%)\]\end{apcanswer}}
  \item \ce{SO2} and \ce{O2} react to completion:
        \ce{2SO2(g) + O2(g) -> 2SO3(g)}. Find the density of the gas
        mixture after the reaction completes (same STP conditions).
        \workspace[6cm]
        \keyonly{\begin{apcanswer}Equimolar \ce{SO2}:\ce{O2} needs a 2:1
        ratio to fully react, so with only 1:1 available,
        \textbf{\ce{SO2} is limiting} --- \ce{O2} is left over.

        \ce{SO3} produced $=m$; \ce{O2} remaining $=m-m/2=m/2$;
        \ce{He} unchanged $=h$.

        New total moles $=m+m/2+h=1.5m+h$. \textbf{Total mass is
        conserved} (no mass enters or leaves):

        \[\text{total mass} = m(64.07)+m(32.00)+h(4.003)\]

        Using the $m/h=3.994$ ratio (take $h=1$, $m=3.994$):

        \[\text{mass}=3.994(96.07)+4.003=387.7\ \text{(arb. units)},
        \quad n_{\text{new}}=1.5(3.994)+1=6.991\]

        \[M_{\text{avg,new}} = \frac{387.7}{6.991} = 55.5\ \text{g/mol}\]

        \[d_{\text{new}} = \frac{M_{\text{avg,new}}}{22.42}
        = \mathbf{2.47}\ \text{g/L}\]

        (Higher than the original \SI{1.924}{\gram\per\liter} --- the
        same total mass is now packed into fewer moles, since
        \ce{2SO2 + O2 -> 2SO3} reduces total gas moles.)
        \end{apcanswer}}
\end{parts}

\begin{rubric}
  \rpoint{1}{(a) finds $M_{\text{avg}}$ of the original mixture from its
    density}
  \rpoint{1}{(a) sets up the mole-fraction equation correctly}
  \rpoint{1}{(a) $\chi(\ce{He})\approx0.111$}
  \rpoint{1}{(b) identifies \ce{SO2} as limiting despite the equimolar
    setup}
  \rpoint{1}{(b) correctly finds the new total moles, $1.5m+h$}
  \rpoint{1}{(b) recognizes total mass is conserved}
  \rpoint{1}{(b) $d_{\text{new}}\approx\SI{2.47}{\gram\per\liter}$}
\end{rubric}

% ---------------------------------------------------------------------------
\frq{long}{7}{Zumdahl \S5.4 \textbullet\ industrial combustion, complete
  and incomplete \emph{(after Zumdahl Q161)}}

\ce{CH4} gas flows into a combustion chamber at \SI{200.}{\liter\per
\minute}, \SI{1.50}{\atm}, \SI{25}{\celsius}. Air (21~mol\% \ce{O2},
79~mol\% \ce{N2}) is added at \SI{1.00}{\atm}, the same temperature, and
the mixture is ignited: \ce{CH4(g) + 2O2(g) -> CO2(g) + 2H2O(g)}.

\begin{parts}
  \item To ensure complete combustion, three times the stoichiometrically
        necessary \ce{O2} is supplied. Find the required air flow rate.
        \workspace[5.5cm]
        \keyonly{\begin{apcanswer}\[\frac{n(\ce{CH4})}{\text{min}}
        = \frac{(1.50)(200.)}{(0.08206)(298)} = 12.27\ \text{mol/min}\]

        Stoichiometric \ce{O2} $=2\times12.27=24.54$~mol/min; actual
        (three times) $=73.61$~mol/min.

        \[\frac{n_{\text{air}}}{\text{min}} = \frac{73.61}{0.21}
        = 350.5\ \text{mol/min}\]

        \[\text{air flow rate} = \frac{nRT}{P}
        = \frac{(350.5)(0.08206)(298)}{1.00}
        = \mathbf{8570}\ \text{L/min}\]\end{apcanswer}}
  \item Under these conditions combustion is incomplete: 95.0\% of the
        exhaust carbon becomes \ce{CO2(g)}, the rest \ce{CO(g)} (via
        \ce{CH4(g) + 1.5O2(g) -> CO(g) + 2H2O(g)}); all \ce{CH4} reacts,
        \ce{N2} is unreacted. Find the mole fraction of \ce{CO},
        \ce{CO2}, \ce{O2}, \ce{N2}, and \ce{H2O} in the exhaust.
        \workspace[7cm]
        \keyonly{\begin{apcanswer}Per minute: \ce{CO2}$=0.95(12.27)
        =11.65$~mol; \ce{CO}$=0.05(12.27)=0.6134$~mol; \ce{H2O}
        $=2(12.27)=24.54$~mol (every \ce{CH4}'s H atoms become \ce{H2O}
        regardless of which carbon product forms).

        \ce{O2} consumed $=0.95(12.27)(2)+0.05(12.27)(1.5)=24.23$~mol/min;
        \ce{O2} remaining $=73.61-24.23=49.38$~mol/min.

        \ce{N2} (unreacted, 79\% of the air found in part a)
        $=0.79\times350.5=276.9$~mol/min.

        Total exhaust $=11.65+0.6134+24.54+49.38+276.9=363.1$~mol/min.

        \[\chi(\ce{CO2})=\mathbf{0.0321},\quad \chi(\ce{CO})=\mathbf{0.0017},
        \quad \chi(\ce{H2O})=\mathbf{0.0676}\]
        \[\chi(\ce{O2})=\mathbf{0.136},\quad
        \chi(\ce{N2})=\mathbf{0.763}\]\end{apcanswer}}
\end{parts}

\begin{rubric}
  \rpoint{1}{(a) moles per minute of \ce{CH4}}
  \rpoint{1}{(a) correctly scales \ce{O2} to 3$\times$ stoichiometric}
  \rpoint{1}{(a) \SI{8570}{\liter\per\minute}}
  \rpoint{1}{(b) correctly splits carbon 95\%/5\% between \ce{CO2} and
    \ce{CO}}
  \rpoint{1}{(b) correctly finds remaining \ce{O2} and unreacted \ce{N2}}
  \rpoint{2}{(b) all five mole fractions correct and summing to 1}
\end{rubric}

% ---------------------------------------------------------------------------
\frq{long}{6}{Zumdahl \S5.4 \textbullet\ graphite combustion, \ce{CO}
  versus \ce{CO2} \emph{(after Zumdahl Q162)}}

A steel cylinder contains \SI{5.00}{\mole} graphite (solid carbon) and
\SI{5.00}{\mole} \ce{O2}. The mixture is ignited; all the graphite
reacts, producing a mixture of \ce{CO} and \ce{CO2} gas. After the
cylinder cools back to its original temperature, its pressure has
increased by 17.0\%.

Calculate the mole fractions of \ce{CO}, \ce{CO2}, and \ce{O2} in the
final gas mixture.
\workspace[7cm]
\keyonly{\begin{apcanswer}Initial \emph{gas} moles $=5.00$ (graphite is
solid). At constant $V,T$, pressure tracks moles, so:

\[n_{\text{final}} = 1.170\times5.00 = 5.85\ \text{mol}\]

Let $x=$~mol \ce{C}$\to$\ce{CO2} (via \ce{C(s) + O2(g) -> CO2(g)}, no net
change in gas moles) and $y=$~mol \ce{C}$\to$\ce{CO} (via
\ce{2C(s) + O2(g) -> 2CO(g)}, a net \emph{gain} of 0.5~mol gas per mol
\ce{C}). Then $x+y=5.00$, and:

\[n_{\text{final}} = \underbrace{[5.00-(x+0.5y)]}_{\ce{O2}\text{ left}}
+\underbrace{x}_{\ce{CO2}}+\underbrace{y}_{\ce{CO}} = 5.00+0.5y\]

\[5.00+0.5y = 5.85 \ \Longrightarrow\ y=\mathbf{1.70}\ \text{mol},
\quad x=5.00-1.70=\mathbf{3.30}\ \text{mol}\]

\[\ce{O2}\ \text{remaining} = 5.00-(3.30+0.85) = 0.85\ \text{mol}\]

\[\chi(\ce{CO2})=\frac{3.30}{5.85}=\mathbf{0.564}, \quad
\chi(\ce{CO})=\frac{1.70}{5.85}=\mathbf{0.291}, \quad
\chi(\ce{O2})=\frac{0.85}{5.85}=\mathbf{0.145}\]\end{apcanswer}}

\begin{rubric}
  \rpoint{1}{recognizes only \ce{O2} counts as initial gas moles}
  \rpoint{1}{correctly finds final gas moles, \SI{5.85}{\mole}, from the
    17.0\% pressure increase}
  \rpoint{2}{sets up and solves $5.00+0.5y=5.85$ for $x$ and $y$}
  \rpoint{1}{correctly finds remaining \ce{O2}}
  \rpoint{1}{all three mole fractions correct and summing to 1}
\end{rubric}

% ---------------------------------------------------------------------------
\frq{long}{7}{Zumdahl \S5.4 \textbullet\ hot-air balloon lift
  \emph{(after Zumdahl Q163)}}

A hot-air balloon approximates a \SI{5.00}{\meter}-diameter sphere,
containing air heated to \SI{65}{\celsius}; surrounding air is
\SI{21}{\celsius}. The balloon's internal pressure equals atmospheric
pressure, \SI{745}{\torr}, throughout. (Average molar mass of air
$=\SI{29.0}{\gram\per\mole}$.)

\begin{parts}
  \item What total mass can the balloon lift?
        \workspace[5.5cm]
        \keyonly{\begin{apcanswer}$V=\tfrac43\pi(2.50)^3=65.45$~m$^3$
        $=65450$~L. $P=745/760=0.9803$~atm.

        \[\text{mass cold air (294 K)} = \frac{PVM}{RT}
        = \frac{(0.9803)(65450)(29.0)}{(0.08206)(294)}
        = 77120\ \text{g}\]

        \[\text{mass hot air (338 K)} = \frac{(0.9803)(65450)(29.0)}
        {(0.08206)(338)} = 67080\ \text{g}\]

        \[\text{lift} = 77120-67080 = \mathbf{1.004\times10^4}\ \text{g}
        \approx\mathbf{10.0}\ \text{kg}\]\end{apcanswer}}
  \item If instead filled with \ce{He} at \SI{21}{\celsius},
        \SI{745}{\torr}, matching the same volume, what total mass can
        the balloon lift?
        \workspace[4.5cm]
        \keyonly{\begin{apcanswer}mass \ce{He} (294~K)
        $=(0.9803)(65450)(4.003)/[(0.08206)(294)]=10650$~g; lift
        $=77120-10650=\mathbf{6.65\times10^4}$~g $\approx\mathbf{66.5}$~kg
        --- over 6 times the hot-air balloon's lift, since \ce{He} is so
        much less dense than even heated ordinary air.\end{apcanswer}}
  \item Repeat part (a) at Denver's typical atmospheric pressure,
        \SI{630.}{\torr}.
        \workspace[4.5cm]
        \keyonly{\begin{apcanswer}$P=630/760=0.8289$~atm; mass cold air
        $=65220$~g, mass hot air $=56730$~g; lift
        $=\mathbf{8.49\times10^3}$~g $\approx\mathbf{8.49}$~kg --- less
        lift than at sea level, since a thinner atmosphere displaces
        less mass to begin with.\end{apcanswer}}
\end{parts}

\begin{rubric}
  \rpoint{1}{correct sphere volume in liters}
  \rpoint{1}{(a) mass of cold (displaced) air}
  \rpoint{1}{(a) mass of hot air inside}
  \rpoint{1}{(a) lift $\approx\SI{10.0}{\kilo\gram}$}
  \rpoint{2}{(b) \ce{He} lift $\approx\SI{66.5}{\kilo\gram}$, correctly
    using \SI{21}{\celsius} for both masses}
  \rpoint{1}{(c) Denver lift $\approx\SI{8.49}{\kilo\gram}$, correctly
    using the lower pressure throughout}
\end{rubric}

% ---------------------------------------------------------------------------
\frq{long}{5}{Zumdahl \S5.4 \textbullet\ balloons needed to lift a child
  \emph{(after Zumdahl Q164)}}

\begin{parts}
  \item Estimate the minimum number of \SI{10.}{\liter} helium balloons
        (filled at \SI{1.00}{\atm}, \SI{25}{\celsius}) needed to lift a
        50.-lb child ($M_{\text{air}}=\SI{29}{\gram\per\mole}$;
        balloon and string masses negligible).
        \workspace[5.5cm]
        \keyonly{\begin{apcanswer}mass air per balloon
        $=PVM/RT=(1.00)(10.0)(29)/[(0.08206)(298)]=11.86$~g

        mass \ce{He} per balloon
        $=(1.00)(10.0)(4.003)/[(0.08206)(298)]=1.637$~g

        lift per balloon $=11.859-1.637=10.222$~g (kept unrounded here so
        the final division comes out right)

        child's mass $=50.\times0.4536=22.68$~kg $=22680$~g

        \[\text{balloons} = \frac{22680}{10.222}
        = 2218.7\ \Longrightarrow\ \text{round up} \Longrightarrow
        \mathbf{2219\ \text{balloons}}\]\end{apcanswer}}
  \item Explain, in terms of buoyancy, why the balloons can lift the
        child at all.
        \workspace[3.5cm]
        \keyonly{\begin{apcanswer}Each balloon displaces a mass of air far greater
        than the mass of helium it contains (\ce{He}'s molar mass is
        about 1/7 that of air). The buoyant force equals the weight of
        the displaced air (Archimedes' principle applied to a gas); the
        excess of that buoyant force over the balloon's own weight is
        what is available to lift additional load.\end{apcanswer}}
\end{parts}

\begin{rubric}
  \rpoint{1}{mass of air displaced per balloon}
  \rpoint{1}{mass of \ce{He} per balloon}
  \rpoint{1}{lift per balloon, $\approx\SI{10.2}{\gram}$}
  \rpoint{1}{correctly converts the child's mass and rounds up to 2219}
  \rpoint{1}{explains the buoyancy argument}
\end{rubric}

% ---------------------------------------------------------------------------
\frq{long}{6}{Zumdahl \S5.4 \textbullet\ matching a hot-air balloon's
  lift to helium's \emph{(after Zumdahl Q165)}}

You have a helium balloon at \SI{1.00}{\atm}, \SI{25}{\celsius}. You want
a hot-air balloon of the \emph{same volume} with the \emph{same} lift.
Air is 79.0\% \ce{N2}, 21.0\% \ce{O2} by mole.

\begin{parts}
  \item Will the hot-air balloon's temperature need to be higher or
        lower than \SI{25}{\celsius}? Explain.
        \workspace[4.5cm]
        \keyonly{\begin{apcanswer}\textbf{Higher.} Ordinary air ($M\approx29$) is
        far denser than \ce{He} ($M\approx4$) at the same conditions. To
        make hot air displace the same net buoyant mass as \ce{He} does,
        its density must be driven down a great deal --- and since
        density $\propto M/T$ at fixed $P$, that means a large increase
        in temperature.\end{apcanswer}}
  \item Calculate the required air temperature.
        \workspace[5.5cm]
        \keyonly{\begin{apcanswer}$M_{\text{air}}=0.79(28.02)+0.21(32.00)
        =28.86$~g/mol.

        Setting the two lifts (mass of displaced ambient air minus mass
        of the balloon's own gas, per unit volume) equal, with the
        displaced air always evaluated at ambient \SI{298}{\kelvin}:

        \[\frac{(M_{\text{air}}-M_{\ce{He}})P}{R(298)}
        = \frac{M_{\text{air}}P}{R(298)} - \frac{M_{\text{air}}P}
        {RT_{\text{hot}}}\]

        $P/R$ cancels completely, leaving
        $M_{\text{air}}/T_{\text{hot}} = M_{\ce{He}}/298$:

        \[T_{\text{hot}} = 298\times\frac{M_{\text{air}}}{M_{\ce{He}}}
        = 298\times\frac{28.86}{4.003}
        = \mathbf{2148}\ \text{K}\ (\mathbf{1875}\ \text{\textcelsius})
        \]

        A physically unreachable temperature for a real balloon
        envelope --- a vivid illustration of just how much less dense
        helium is than even very hot air.\end{apcanswer}}
\end{parts}

\begin{rubric}
  \rpoint{1}{(a) correctly identifies "higher," with density reasoning}
  \rpoint{1}{(b) average molar mass of air}
  \rpoint{2}{(b) correctly sets the two lift expressions equal}
  \rpoint{2}{(b) $T_{\text{hot}}\approx\SI{2148}{\kelvin}$, algebra shown}
\end{rubric}

% ---------------------------------------------------------------------------
\frq{long}{4}{Zumdahl \S5.8 \textbullet\ van der Waals reduces to the
  ideal gas law \textbf{[beyond CED]} \emph{(after Zumdahl Q166)}}

\[\left[P+a\left(\frac{n}{V}\right)^2\right](V-nb) = nRT\]

Show algebraically that this reduces to $PV=nRT$ at low pressure (large
$V$, fixed $n$) and explain why the same is true at high temperature.
\workspace[6.5cm]
\keyonly{\begin{apcanswer}\textbf{Low pressure / large volume.} As $V$
grows without bound at fixed $n$: the attraction term $a(n/V)^2\to0$,
since $n/V\to0$; and $nb$ becomes negligible \emph{relative to} $V$, so
$(V-nb)\to V$. Substituting both limits:

\[[P+0](V-0) \ \longrightarrow\ PV=nRT\]

Physically: at low pressure molecules are far apart on average, so (i)
their own volume $nb$ is a negligible fraction of the container, and
(ii) intermolecular attractions --- which fall off steeply with distance
--- rarely have a chance to act.

\textbf{High temperature.} The attraction term $a(n/V)^2$ does not
depend on $T$ at all, but $nRT$ grows without bound as $T$ rises, so the
attraction term's contribution to the \emph{measured} pressure becomes
comparatively negligible: kinetic energy increasingly dominates over the
potential energy of intermolecular attraction. (The $nb$ term still
requires $V$ to stay reasonably large --- i.e., pressure not also pushed
to an extreme --- for it to vanish too.) Both limits describe the same
underlying idea: the two correction terms exist specifically to capture
molecular size and intermolecular attraction, and both become
unimportant whenever the gas is dilute or energetic enough that
molecules spend most of their time far apart and moving fast relative to
any attraction between them.\end{apcanswer}}

\begin{rubric}
  \rpoint{1}{shows $a(n/V)^2\to0$ as $V\to\infty$}
  \rpoint{1}{shows $(V-nb)\to V$ as $V\to\infty$, reaching $PV=nRT$}
  \rpoint{1}{explains the low-pressure limit physically (molecules far
    apart, both corrections vanish)}
  \rpoint{1}{explains why high $T$ also drives the equation toward ideal
    behavior}
\end{rubric}

% ---------------------------------------------------------------------------
\frq{long}{5}{Zumdahl \S5.4 \textbullet\ identifying an unknown compound
  \emph{(after Zumdahl Q167)}}

An unknown gaseous binary compound (two elements) is burned in excess
oxygen. Burning \SI{10.0}{\gram} of it produces \SI{16.3}{\gram} water.
The compound's density is 1.38 times that of oxygen gas at the same
temperature and pressure.

Give a possible identity for the unknown compound.
\workspace[6.5cm]
\keyonly{\begin{apcanswer}\[M_{\text{compound}} = 1.38\times32.00
= 44.16\ \text{g/mol}\]

Since the compound produces water on combustion, it must contain
\textbf{hydrogen} --- this immediately rules out close-mass decoys like
\ce{N2O} and \ce{CO2} (both $\approx44.01$~g/mol, neither containing H).

\textbf{Test propane, \ce{C3H8}} ($M=3(12.01)+8(1.008)=44.09$~g/mol,
within 0.15\% of the density-derived value):

\[\ce{C3H8(g) + 5O2(g) -> 3CO2(g) + 4H2O(g)}\]

\[n(\ce{C3H8}) = \frac{10.0}{44.09} = 0.2268\ \text{mol}\]

\[\text{predicted mass \ce{H2O}} = 4(0.2268)(18.02) = 16.3\ \text{g}\]

This matches the given \SI{16.3}{\gram} almost exactly.

\[\textbf{Identity: \ce{C3H8} (propane)}\]\end{apcanswer}}

\begin{rubric}
  \rpoint{1}{molar mass of the unknown, $\approx\SI{44.16}{\gram\per
    \mole}$}
  \rpoint{1}{recognizes the compound must contain H, ruling out
    mass-matched non-hydrogen decoys}
  \rpoint{1}{balances the combustion equation for \ce{C3H8}}
  \rpoint{1}{predicts the water mass from the combustion stoichiometry}
  \rpoint{1}{identifies \ce{C3H8} and confirms the match}
\end{rubric}

% ---------------------------------------------------------------------------
\frq{long}{6}{Zumdahl \S5.4 \textbullet\ ammonia synthesis at constant
  pressure \emph{(after Zumdahl Q168)}}

\ce{N2(g) + 3H2(g) -> 2NH3(g)}. \SI{15.0}{\liter} of \ce{N2} and \ce{H2}
gas sit in a container fitted with a movable piston that holds the total
pressure constant. Each reactant initially has a partial pressure of
\SI{1.00}{\atm}.

\begin{parts}
  \item Find the partial pressure of \ce{NH3} after the reaction
        completes.
        \workspace[6cm]
        \keyonly{\begin{apcanswer}Equal initial partial pressures at the
        same $V,T$ means equal initial moles, $n_0$ each. The 3:1
        stoichiometric requirement means \textbf{\ce{H2} is limiting}
        (only enough to react $n_0/3$~mol \ce{N2}).

        All $n_0$~mol \ce{H2} reacts, consuming $n_0/3$~mol \ce{N2} and
        producing $2n_0/3$~mol \ce{NH3}. \ce{N2} remaining
        $=n_0-n_0/3=2n_0/3$.

        Final total moles $=2n_0/3\ (\ce{N2}) + 2n_0/3\ (\ce{NH3})
        = 4n_0/3$ (vs.\ $2n_0$ initially).

        Total pressure stays fixed at its initial value,
        $1.00+1.00=\SI{2.00}{\atm}$ (the piston's job). \ce{NH3}'s mole
        fraction in the final mixture is $(2n_0/3)/(4n_0/3)=0.500$:

        \[P(\ce{NH3}) = 2.00\times0.500 = \mathbf{1.00}\ \text{atm}\]
        \end{apcanswer}}
  \item Find the volume of the container after the reaction completes.
        \workspace[3.5cm]
        \keyonly{\begin{apcanswer}At constant $T,P$, $V\propto n$:

        \[V_{\text{final}} = 15.0\times\frac{4n_0/3}{2n_0}
        = 15.0\times\frac23 = \mathbf{10.0}\ \text{L}\]\end{apcanswer}}
\end{parts}

\begin{rubric}
  \rpoint{1}{identifies \ce{H2} as limiting from the 3:1 stoichiometric
    requirement against equal initial moles}
  \rpoint{1}{correctly finds the final moles of each species in terms of
    $n_0$}
  \rpoint{1}{finds the final total moles, $4n_0/3$}
  \rpoint{1}{(a) $P(\ce{NH3})=\SI{1.00}{\atm}$, via the mole-fraction
    argument}
  \rpoint{2}{(b) $V_{\text{final}}=\SI{10.0}{\liter}$}
\end{rubric}

% ---------------------------------------------------------------------------
\frq{long}{4}{Zumdahl \S5.7 \textbullet\ diffusion meeting point
  \textbf{[beyond CED]} \emph{(after Zumdahl Q169)}}

A classroom has ten evenly spaced rows of students. A student in row~1
releases \ce{N2O} (laughing gas) and a student in row~10 simultaneously
releases a lachrymator (tear-inducing gas) of molar mass
\SI{176}{\gram\per\mole}.

In which row do students first laugh and cry at the same time?
\workspace[5cm]
\keyonly{\begin{apcanswer}Diffusion distance is proportional to
$1/\sqrt M$, so the two gases split the 9 row-gaps between rows 1 and 10
in the ratio of their diffusion rates:

\[\frac{d_{\ce{N2O}}}{d_{\text{lachrymator}}}
= \sqrt{\frac{M_{\text{lachrymator}}}{M_{\ce{N2O}}}}
= \sqrt{\frac{176}{44.01}} = 2.00\]

With $d_{\ce{N2O}}+d_{\text{lachrymator}}=9$ and
$d_{\ce{N2O}}=2\,d_{\text{lachrymator}}$:

\[d_{\ce{N2O}} = 9\times\frac{2}{3} = 6.00\ \text{row-gaps}\]

\[\textbf{Meeting row} = 1+6.00 = \mathbf{\text{row } 7}\]\end{apcanswer}}

\begin{rubric}
  \rpoint{1}{correct Graham's-law-based distance ratio,
    $\sqrt{176/44.01}\approx2$}
  \rpoint{1}{sets up $d_{\ce{N2O}}+d_{\text{lachrymator}}=9$}
  \rpoint{1}{solves for $d_{\ce{N2O}}=6$ row-gaps}
  \rpoint{1}{correctly reports row 7}
\end{rubric}

% ============================================================================
\examtotal

\end{document}
